\subsection{Local cohomology}

In this number, we consider a local Noetherian ring A. Recall (VIII, § 3, n° 3, lemma 2) that the ideals of definition of A are the ideals of A distinct from A containing a power of \(\mathfrak{m}_{\mathrm{A}}\) , or equivalently the ideals \(\mathfrak{a} \subset \mathfrak{m}_{\mathrm{A}}\) such that \(\mathrm{A}/\mathfrak{a}\) is of finite length. We will denote \(\mathcal{D}\) the set of defining ideals of \(\mathrm{A}\) , equipped with the order relation opposite to inclusion; it is filtered on the right.

Let M be an A-module. Associate with every ideal of definition a of A the graded A-module \(\operatorname{Ext}_{\mathrm{A}}(\mathrm{A}/\mathfrak{a},\mathrm{M})\) ; if a and b are ideals of definition with \(a \subset b\) , let us denote by \(p_{ab}: A/a \to A/b\) the canonical map, and consider the A-linear map \(\operatorname{Ext}(p_{ab},1_{\mathrm{M}}): \operatorname{Ext}_{\mathrm{A}}(\mathrm{A}/\mathfrak{b},\mathrm{M}) \longrightarrow \operatorname{Ext}_{\mathrm{A}}(\mathrm{A}/\mathfrak{a},\mathrm{M})\) . We thus obtain an inductive system of graded A-modules and graded A-linear maps of degree 0 relative to the ordered set D. We call the local cohomology module of M, and denote \(\mathrm{H}_{\mathrm{A}}(\mathrm{M})\) , the graded A-module \(\lim \operatorname{Ext}_{\mathrm{A}}(\mathrm{A}/\mathfrak{a},\mathrm{M})\) .

The ideals \(\mathfrak{m}_{\mathrm{A}}^{n}\) for \(n \geqslant 1\) form a cofinal subset of \(\mathcal{D}\) ; we therefore have a canonical isomorphism of graded modules \(\varinjlim_{n} \operatorname{Ext}_{\mathrm{A}}(\mathrm{A}/\mathfrak{m}_{\mathrm{A}}^{n}, M) \longrightarrow \mathrm{H}_{\mathrm{A}}(\mathrm{M})\) . Consequently, every element of \(\mathrm{H}_{\mathrm{A}}(\mathrm{M})\) is annihilated by a power of \(\mathfrak{m}_{\mathrm{A}}\) .

Remark 1.--- *Let X be the topological space* \(\operatorname{Spec}(A)\) , \(\mathcal{O}_{\mathrm{X}}\) the structural sheaf of rings and \(\widetilde{\mathbf{M}}\) the \(\mathcal{O}_{\mathrm{X}}\) -module associated with M. The graded A-module \(\mathrm{H}_{\mathrm{A}}(\mathrm{M})\) is identified with the module \(\mathrm{H}_{\{\mathfrak{m}_{\mathrm{A}}\}}(\mathrm{X},\widetilde{\mathrm{M}})\) of cohomology with support in the closed point \(\mathfrak{m}_{\mathrm{A}}\) of \(X\) .

For every homomorphism \(f: \mathrm{M} \to \mathrm{N}\) of A-modules, the maps \(\operatorname{Ext}_{\mathrm{A}}(1_{\mathrm{A}/\mathfrak{a}}, f): \operatorname{Ext}_{\mathrm{A}}(\mathrm{A}/\mathfrak{a}, \mathrm{M}) \longrightarrow \operatorname{Ext}_{\mathrm{A}}(\mathrm{A}/\mathfrak{a}, \mathrm{N})\) form an inductive system of graded linear maps. Passing to the inductive limit, one obtains a graded homomorphism \(\mathrm{H}_{\mathrm{A}}(f): \mathrm{H}_{\mathrm{A}}(\mathrm{M}) \to \mathrm{H}_{\mathrm{A}}(\mathrm{N})\) . For every sequence \(\mathrm{M} \xrightarrow{f} \mathrm{N} \xrightarrow{g} \mathrm{P}\) of A-modules and homomorphisms, one has \(\mathrm{H}_{\mathrm{A}}(g \circ f) = \mathrm{H}_{\mathrm{A}}(g) \circ \mathrm{H}_{\mathrm{A}}(f)\) . Let

\[
0 \longrightarrow \mathrm{M} \stackrel {f} {\longrightarrow} \mathrm{N} \stackrel {g} {\longrightarrow} \mathrm{P} \longrightarrow 0\tag{ℰ}
\]

an exact sequence of A-modules. By A, X, p.90, prop. 8, the connecting homomorphisms of the extension modules \(\operatorname{Ext}_{\mathrm{A}}(\mathrm{A}/\mathfrak{a},\mathrm{P}) \longrightarrow \operatorname{Ext}_{\mathrm{A}}(\mathrm{A}/\mathfrak{a},\mathrm{M})\) form an inductive system of A-linear maps, graded of (ascending) degree +1. Passing to the inductive limit, one deduces an A-homomorphism \(\partial(\mathcal{E}): \mathrm{H}_{\mathrm{A}}(\mathrm{P}) \to \mathrm{H}_{\mathrm{A}}(\mathrm{M})\) , graded of degree +1, which makes exact the sequence of homomorphisms

\[
\dots \longrightarrow \mathrm{H} _ {\mathrm{A}} ^ {n - 1} (\mathrm{P}) \xrightarrow {\partial^ {n - 1} (\mathcal {E})} \mathrm{H} _ {\mathrm{A}} ^ {n} (\mathrm{M}) \xrightarrow {\mathrm{H} _ {\mathrm{A}} ^ {n} (f)} \mathrm{H} _ {\mathrm{A}} ^ {n} (\mathrm{N}) \xrightarrow {\mathrm{H} _ {\mathrm{A}} ^ {n} (g)} \mathrm{H} _ {\mathrm{A}} ^ {n} (\mathrm{P}) \xrightarrow {\partial^ {n} (\mathcal {E})} \mathrm{H} _ {\mathrm{A}} ^ {n + 1} (\mathrm{M}) \longrightarrow \dots
\]

Let M be an A-module. For every ideal a of A, the A-module \(\mathrm{Hom}_{\mathrm{A}}(\mathrm{A}/\mathfrak{a},\mathrm{M})\) is canonically identified with the submodule of M consisting of the elements annihilated by a. Thus \(\mathrm{H}_{\mathrm{A}}^{0}(\mathrm{M})\) identifies with the submodule of M consisting of the elements m that are annihilated by a power of \(\mathfrak{m}_{\mathrm{A}}\) , that is, such that \(\mathrm{long}_{\mathrm{A}}(\mathrm{Am}) < +\infty\) . In particular, we have \(\mathrm{H}_{\mathrm{A}}^{0}(\mathrm{M}) = \mathrm{M}\) when M is Artinian.

Examples.--- 1) If M is injective, the A-module \(\mathrm{H}_{\mathrm{A}}^{i}(\mathrm{M})\) is zero for \(i > 0\) and injective for \(i = 0\) ( \(\S 8, \mathrm{n}^{\circ}2\) , lemma 1, c)).

\begin{enumerate}
\def\labelenumi{\arabic{enumi})}
\setcounter{enumi}{1}
\item
  If \(\mathrm{H}_{\mathrm{A}}^{0}(\mathrm{M})=\mathrm{M}\) (for example if M is Artinian), \(\mathrm{H}_{\mathrm{A}}^{i}(\mathrm{M})\) is zero for i > 0. Let (I, e) be an injective envelope of M. The submodule \(\mathrm{H}_{\mathrm{A}}^{0}(\mathrm{I})\) of I is injective (example 1) and contains \(e(\mathrm{M})\) , therefore is equal to I. Set N to be the cokernel of e and consider the exact sequence \(0\rightarrow M\xrightarrow{e}\mathrm{I}\xrightarrow{p}\mathrm{N}\rightarrow0\) . Since \(\mathrm{I}=\mathrm{H}_{\mathrm{A}}^{0}(\mathrm{I})\) , one has \(\mathrm{N}=\mathrm{H}_{\mathrm{A}}^{0}(\mathrm{N})\) and the homomorphism \(\mathrm{H}_{\mathrm{A}}^{0}(p)\) is surjective. Since \(\mathrm{H}_{\mathrm{A}}^{i}(\mathrm{I})\) is zero for i > 0 (example 1), \(\mathrm{H}_{\mathrm{A}}^{1}(\mathrm{M})\) is zero, and \(\mathrm{H}_{\mathrm{A}}^{i}(\mathrm{M})\) is isomorphic to \(\mathrm{H}_{\mathrm{A}}^{i-1}(\mathrm{N})\) for i > 1; we conclude by reasoning by induction on the integer i.
\item
  Let \(\Omega\) be a dualizing A-module. For \(i \neq \dim(A)\) , one has \(\operatorname{Ext}_{A}^{i}(A/a, \Omega) = 0\) for every ideal of definition \(a\) of A ( \(\S 8\) , \(n^{\circ} 5\) , example 3), whence \(\mathrm{H}_{\mathrm{A}}^{i}(\Omega) = 0\) ; for \(i = \dim(A)\) , the A-module \(\mathrm{H}_{\mathrm{A}}^{i}(\Omega)\) , which is isomorphic to \(\varinjlim \operatorname{Ext}_{\mathrm{A}}^{i}(\mathrm{A}/\mathfrak{m}_{\mathrm{A}}^{n}, \Omega)\) , is a Matlis A-module (loc. cit., example 6).
\item
  Let A be a Noetherian local integral domain; denote by K its field of fractions, and suppose A ≠ K. K is an injective A-module (A, X, p.18, example 1), so that the module \(\mathrm{H}_{\mathrm{A}}(K)\) is zero (example 1). From the exact sequence 0 → A → K → K/A → 0, we derive for every i an isomorphism \(\mathrm{H}_{\mathrm{A}}^{i}(K/A)\) → \(\mathrm{H}_{\mathrm{A}}^{i+1}(A)\).
\end{enumerate}

More generally, for every torsion-free A-module M and every integer i, we deduce from the exact sequence

\[
0 \rightarrow \mathrm{M} \rightarrow \mathrm{K} \otimes_ {\mathrm{A}} \mathrm{M} \rightarrow (\mathrm{K/A}) \otimes_ {\mathrm{A}} \mathrm{M} \rightarrow 0
\]

an isomorphism \(\mathrm{H}_{\mathrm{A}}^{i}((\mathrm{K}/\mathrm{A})\otimes_{\mathrm{A}}\mathrm{M})\to \mathrm{H}_{\mathrm{A}}^{i + 1}(\mathrm{M})\)

\begin{enumerate}
\def\labelenumi{\arabic{enumi})}
\setcounter{enumi}{4}
\tightlist
\item
  Keep the hypotheses of the preceding example and suppose moreover that dim(A) = 1. Let N be a torsion A-module; since every nonzero ideal of A distinct from A is an ideal of definition (VIII, § 1, n° 3, prop. 6, e)), we have \(\mathrm{H}_{\mathrm{A}}^{0}(N) = N\), and consequently \(\mathrm{H}_{\mathrm{A}}^{i}(N) = 0\) for i > 0 (example 2).
\end{enumerate}

Let M be an A-module; denote by T(M) its torsion submodule. Consider the long exact sequence of local cohomology associated with the exact sequence

\[
0 \to \mathrm{T} (\mathrm{M}) \to \mathrm{M} \to \mathrm{M} / \mathrm{T} (\mathrm{M}) \to 0;
\]

taking the preceding into account, we deduce from it canonical isomorphisms T(M) → \(\mathrm{H}_{\mathrm{A}}^{0}(M)\) and \(\mathrm{H}_{\mathrm{A}}^{1}(M)\) → \(\mathrm{H}_{\mathrm{A}}^{1}(M/T(M))\). Since the canonical homomorphism \((K/A) \otimes_{A} M\) → \((K/A) \otimes_{A} (M/T(M))\) is bijective, we finally obtain canonical isomorphisms

\[
\mathrm{H} _ {\mathrm{A}} ^ {0} (\mathrm{M}) \rightarrow \mathrm{T} (\mathrm{M}), \quad \mathrm{H} _ {\mathrm{A}} ^ {1} (\mathrm{M}) \rightarrow (\mathrm{K} / \mathrm{A}) \otimes_ {\mathrm{A}} \mathrm{M}.
\]

PROPOSITION 1. Let A be a Noetherian local ring and M a finitely generated A-module.

\begin{enumerate}
\def\labelenumi{\alph{enumi})}
\item
  The A-module \(\mathrm{H}_{\mathrm{A}}(\mathrm{M})\) is Artinian, and zero in degree \(>\dim (\mathrm{M})\) .
\item
  Set \(p = \mathrm{prof}_{\mathrm{A}}(\mathrm{M})\) . One has \(\mathrm{H}_{\mathrm{A}}^{i}(\mathrm{M}) = 0\) for \(i < p\) , and \(\mathrm{H}_{\mathrm{A}}^{p}(\mathrm{M}) \neq 0\) if \(M\) is nonzero.
\end{enumerate}

Let us prove a) by reasoning by induction on \(\dim(M)\) . The case \(\dim(M) \leqslant 0\) follows from example 2 above. Suppose \(\dim(M) > 0\) and first take \(M\) of the form \(A/p\) , where \(p\) is a prime ideal of \(A\) distinct from \(\mathfrak{m}_A\) . Let \(x\) be an element of \(\mathfrak{m}_A - p\) ; we have an exact sequence \(0 \to M \xrightarrow{x_M} M \longrightarrow M/xM \to 0\) , with \(\dim(M/xM) = \dim(M) - 1\) . From this we deduce an exact sequence of local cohomology

\[
\mathrm{H} _ {\mathrm{A}} ^ {i - 1} (\mathrm{M} / x \mathrm{M}) \longrightarrow \mathrm{H} _ {\mathrm{A}} ^ {i} (\mathrm{M}) \stackrel {x} {\longrightarrow} \mathrm{H} _ {\mathrm{A}} ^ {i} (\mathrm{M}).
\]

Every element of \(\mathrm{H}_{\mathrm{A}}^{i}(\mathrm{M})\) is annihilated by a power of \(\mathfrak{m}_{\mathrm{A}}\) ; to prove that this module is Artinian, it therefore suffices to prove that the socle of \(\mathrm{H}_{\mathrm{A}}^{i}(\mathrm{M})\) is of finite dimension over \(\kappa_{A}\) (§ 8, n° 3, lemma 3). By the induction hypothesis, the kernel N of the homothety of ratio x in \(\mathrm{H}_{\mathrm{A}}^{i}(\mathrm{M})\) is Artinian; since x belongs to \(\mathfrak{m}_{\mathrm{A}}\) , the socle of \(\mathrm{H}_{\mathrm{A}}^{i}(\mathrm{M})\) identifies with that of N, hence is of finite dimension. If i > dim(M), we have \(\mathrm{H}_{\mathrm{A}}^{i-1}(\mathrm{M}/x\mathrm{M}) = 0\) by the induction hypothesis, so that the homothety of ratio x is injective in \(\mathrm{H}_{\mathrm{A}}^{i}(\mathrm{M})\) ; since every element of \(\mathrm{H}_{\mathrm{A}}^{i}(\mathrm{M})\) is annihilated by a power of x, we deduce \(\mathrm{H}_{\mathrm{A}}^{i}(\mathrm{M}) = 0\) , whence a) in the case under consideration.

Let us pass to the general case. The A-module M admits a composition series \((\mathrm{M}_{j})_{0\leqslant j\leqslant n}\) such that each quotient \(M_{j}/M_{j+1}\) is isomorphic to \(A/p_{j}\) , where \(p_{j}\) is a prime ideal of A (IV, § 1, n° 4, th. 1). Let us prove by induction on n that M satisfies a). The case n=0 is trivial. The exact sequence \(0\to M_{1}\to M\to A/p_{0}\to0\) provides an exact sequence of local cohomology

\[
\mathrm{H} _ {\mathrm{A}} ^ {i} \left(\mathrm{M} _ {1}\right) \longrightarrow \mathrm{H} _ {\mathrm{A}} ^ {i} (\mathrm{M}) \longrightarrow \mathrm{H} _ {\mathrm{A}} ^ {i} \left(\mathrm{A} / \mathfrak {p} _ {0}\right);
\]

the A-module \(\mathrm{H}_{\mathrm{A}}^{i}(\mathrm{M}_{1})\) is Artinian by the induction hypothesis, and the same holds for \(\mathrm{H}_{\mathrm{A}}^{i}(\mathrm{A}/\mathfrak{p}_{0})\) by the cases already treated; consequently \(\mathrm{H}_{\mathrm{A}}^{i}(\mathrm{M})\) is Artinian. If \(i > \dim(\mathrm{M})\) the modules \(M_{1}\) and \(A/p_{0}\) are of dimension < i; the modules \(\mathrm{H}_{\mathrm{A}}^{i}(\mathrm{M}_{1})\) and \(\mathrm{H}_{\mathrm{A}}^{i}(\mathrm{A}/\mathfrak{p}_{0})\) are therefore zero by the induction hypothesis and the cases already treated, which implies \(\mathrm{H}_{\mathrm{A}}^{i}(\mathrm{M}) = 0\) .

Suppose M is nonzero, and prove b) by induction on the integer \(p = \text{prof}(M)\) . The case \(p = 0\) follows from the definition of depth. Suppose \(p > 0\) and let us choose an element \(x\) of \(\mathfrak{m}_{\mathrm{A}}\) such that the homothety \(x_{\mathrm{M}}\) is injective. As above, we obtain an exact sequence of local cohomology.

\[
\mathrm{H} _ {\mathrm{A}} ^ {i - 1} (\mathrm{M} / x \mathrm{M}) \longrightarrow \mathrm{H} _ {\mathrm{A}} ^ {i} (\mathrm{M}) \xrightarrow {x} \mathrm{H} _ {\mathrm{A}} ^ {i} (\mathrm{M}).
\]

We have \(\operatorname{prof}(M/xM) = \operatorname{prof}(M) - 1\) ( \(\S 1, n^{\circ} 4\) prop. 7), whence \(H_{A}^{i-1}(M/xM) = 0\) for \(i < p\) by the induction hypothesis, which implies as above \(H_{A}^{i}(M) = 0\) . In particular \(H_{A}^{p-1}(M)\) is zero, so that the homomorphism \(H_{A}^{p-1}(M/xM) \longrightarrow H_{A}^{p}(M)\) is injective; thus \(H_{A}^{p}(M)\) is nonzero by the induction hypothesis.

One can show that the module \(H_{A}^{\dim(M)}(M)\) is nonzero when M is nonzero (exerc. 4; cf. no. 3, cor. of th. 2).

COROLLARY. Let M be a nonzero finitely generated Macaulay A-module. The A-module \(\mathrm{H}_{\mathrm{A}}^{i}(\mathrm{M})\) is zero for \(i \neq \dim(\mathrm{M})\) and nonzero for \(i = \dim(\mathrm{M})\) .

Remark 2.--- For every ideal of definition a of A, the A-module \(\operatorname{Ext}_{\mathbf{A}}(\mathbf{A}/\mathfrak{a},\mathbf{M})\) is annihilated by a, and A/a is identified with \(\widehat{\mathbf{A}}/\mathfrak{a}\widehat{\mathbf{A}}\) ; therefore, the graded A-module \(\operatorname{Ext}_{\mathbf{A}}(\mathbf{A}/\mathfrak{a},\mathbf{M})\) is identified with \(\widehat{\mathbf{A}}\otimes_{\mathbf{A}}\operatorname{Ext}_{\mathbf{A}}(\mathbf{A}/\mathfrak{a},\mathbf{M})\) , and hence also with \(\operatorname{Ext}_{\widehat{\mathbf{A}}}(\widehat{\mathbf{A}}/\mathfrak{a}\widehat{\mathbf{A}},\widehat{\mathrm{A}}\otimes_{\mathbf{A}}\mathbf{M})\) (A, X, p.111, prop. 10). The set of ideals \(\mathfrak{a}\widehat{\mathbf{A}}\) , for \(\mathfrak{a}\in\mathcal{D}\) , contains the powers of \(\mathfrak{m}_{\widehat{\mathbf{A}}}\) , hence is cofinal in the set of defining ideals of \(\widehat{\mathbf{A}}\) ; we therefore deduce from the preceding an isomorphism of graded A-modules

\[
\mathrm{H} _ {\mathrm{A}} (\mathrm{M}) \longrightarrow \mathrm{H} _ {\widehat {\mathrm{A}}} (\widehat {\mathrm{A}} \otimes_ {\mathrm{A}} \mathrm{M}).
\]

If the A-module M is of finite type, the A-module \(\widehat{\mathrm{A}}\otimes_{\mathrm{A}}\mathrm{M}\) is identified with the completion \(\widehat{\mathrm{M}}\) of M (III, § 3, n° 4, th. 3), and we have an isomorphism \(\mathrm{H}_{\mathrm{A}}(\mathrm{M})\to \mathrm{H}_{\widehat{\mathrm{A}}}(\widehat{\mathrm{M}})\) , graded of degree 0.

\subsection{Local cohomology over a Macaulay ring}

In this section, we assume that A is a local Macaulay ring; we set \(\dim(A) = d\) .

The ideals generated by a sequence of elements of \(\mathfrak{m}_{\mathrm{A}}\) completely secant for A and of length \(d = \dim (\mathrm{A})\) form a cofinal subset \(\mathcal{D}_{cs}\) in the set \(\mathcal{D}\) of defining ideals of A. Indeed, let \((x_{1},\ldots ,x_{d})\) be a sequence of elements of \(\mathfrak{m}_{\mathrm{A}}\) completely secant for A (§ 2, n° 3, prop. 3); for every integer \(n\) , the sequence \((x_1^n,\dots ,x_d^n)\) is completely secant for A (A, X, p.158, prop. 6, c)), and generates a defining ideal (VIII, § 3, n° 2, cor. of prop. 3 and th. 1) contained in \(\mathfrak{m}_{\mathrm{A}}^n\) .

Let \(\mathfrak{a} \in \mathcal{D}_{cs}\) , and let \(\pi: L \to A/\mathfrak{a}\) be a free resolution of finite type, zero in degree \(>d\) (for example the Koszul complex associated with a completely secant sequence for A generating \(\mathfrak{a}\) ). Consider the dual \(L^{*} = \operatorname{Homgr}_{A}(L,A)\) of \(L\) ; since the depth of A is equal to \(d\) , one has \(\operatorname{Ext}_{A}^{i}(A/\mathfrak{a}, A) = 0\) for \(i < d\) ( \(\S\) 1, \(n^{\circ}\) 1, cor. 2 of prop. 2). Since \(L^{*}\) is of length \(\leqslant d\) , we deduce that \(H^{i}(L^{*})\) is zero for \(i \neq d\) and that \(H^{d}(L^{*})\) is identified with \(\operatorname{Ext}_{A}^{d}(A/\mathfrak{a}, A)\) (A, X, p.100, th. 1). Consequently, we have a homomorphism

\[
\pi^ {*}: \mathrm{L} ^ {*} (- d) \longrightarrow \operatorname{Ext} _ {\mathrm{A}} ^ {d} (\mathrm{A} / \mathfrak {a}, \mathrm{A})
\]

which defines a free resolution of finite type of \(\operatorname{Ext}_{\mathrm{A}}^{d}(\mathrm{A}/\mathfrak{a},\mathrm{A})\) .

Let M be an A-module; consider the canonical isomorphisms (loc. cit.)

\[
\varphi (\mathrm{L}, \mathrm{M}): \mathrm{H} (\operatorname{Homgr} _ {\mathrm{A}} (\mathrm{L}, \mathrm{M})) \rightarrow \operatorname{Ext} _ {\mathrm{A}} (\mathrm{A} / \mathfrak {a}, \mathrm{M})
\]

\[
\psi (\mathrm{M}, \mathrm{L} ^ {*} (- d)): \mathrm{Tor} ^ {\mathrm{A}} (\mathrm{M}, \mathrm{Ext} _ {\mathrm{A}} ^ {d} (\mathrm{A} / \mathfrak {a}, \mathrm{A})) \to \mathrm{H} (\mathrm{M} \otimes_ {\mathrm{A}} \mathrm{L} ^ {*}) (- d).
\]

Since the complex L is free of finite type, the canonical morphism of complexes \(M \otimes_{A} L^{*} \to \text{Homgr}_{A}(L, M)\) is an isomorphism; from this we deduce an isomorphism of graded A-modules \(H(M \otimes_{A} L^{*}) \to H(\text{Homgr}_{A}(L, M))\) . By composing the preceding isomorphisms, we obtain an isomorphism of graded A-modules, said to be canonical.

\[
\tau (\mathrm{L}, \mathrm{M}): \operatorname{Tor} ^ {\mathrm{A}} (\mathrm{M}, \operatorname{Ext} _ {\mathrm{A}} ^ {d} (\mathrm{A} / \mathfrak {a}, \mathrm{A})) (d) \longrightarrow \operatorname{Ext} _ {\mathrm{A}} (\mathrm{A} / \mathfrak {a}, \mathrm{M})
\]

which induces, for each integer i, an isomorphism

\[
\tau^ {i} (\mathrm{L}, \mathrm{M}): \operatorname{Tor} _ {d - i} ^ {\mathrm{A}} (\mathrm{M}, \operatorname{Ext} _ {\mathrm{A}} ^ {d} (\mathrm{A} / \mathfrak {a}, \mathrm{A})) \longrightarrow \operatorname{Ext} _ {\mathrm{A}} ^ {i} (\mathrm{A} / \mathfrak {a}, \mathrm{M}).
\]

For \(M = A\) , \(\tau^d(L, A)\) is the canonical isomorphism from \(A \otimes_A \operatorname{Ext}_A^d(\mathrm{A} / \mathfrak{a}, A)\) onto \(\operatorname{Ext}_\mathrm{A}^d(A / \mathfrak{a}, A)\) .

Let \(\mathfrak{b}\) be an ideal of \(\mathcal{D}_{cs}\) contained in \(\mathfrak{a}\) . Let \(\rho: R \to A/\mathfrak{b}\) be a finite free resolution of length \(\leqslant d\) and let \(p_{\mathfrak{ab}}: A/\mathfrak{b} \to A/\mathfrak{a}\) be the canonical surjection. By A, X, p.49, prop. 3, there exists a morphism of complexes \(\mathrm{P}_{\mathrm{LR}}: R \to L\) such that \(\pi \circ \mathrm{P}_{\mathrm{LR}} = p_{\mathfrak{ab}} \circ \rho\) . By prop. 2 of A, X, p.103, we have a commutative diagram

\[
\begin{array}{c c c} \operatorname{Tor} ^ {A} (M, \operatorname{Ext} _ {A} ^ {d} (A / \mathfrak {a}, A)) (d) & \xrightarrow {\operatorname{Tor} (1 _ {M} , \operatorname{Ext} ^ {d} (p _ {a b} , 1 _ {\mathrm{A}}))} & \operatorname{Tor} ^ {\mathrm{A}} (M, \operatorname{Ext} _ {A} ^ {d} (A / \mathfrak {b}, A)) (d) \\ \psi (L ^ {*} (- d), M) \Bigg \downarrow & & \Bigg \downarrow \psi (R ^ {*} (- d), M) \\ H (M \otimes_ {A} L ^ {*}) & \xrightarrow {H (1 _ {M} \otimes^ {t} P _ {L R})} & H (M \otimes_ {A} R ^ {*}) \\ \Bigg \downarrow & & \Bigg \downarrow \\ H (\operatorname{Homgr} _ {\mathrm{A}} (L, M)) & \xrightarrow {H (\operatorname{Homgr} (P _ {L R} , M))} & H (\operatorname{Homgr} _ {A} (R, M)) \\ \varphi (L, M) \Bigg \downarrow & & \Bigg \downarrow \varphi (R, M) \\ \operatorname{Ext} _ {\mathrm{A}} (A / \mathfrak {a}, M) & \xrightarrow {\operatorname{Ext} (p _ {a b} , 1 _ {M})} & \operatorname{Ext} _ {A} (A / \mathfrak {b}, M) \end{array}
\]

It follows first, by taking \(\mathfrak{a} = \mathfrak{b}\) , that the isomorphism \(\tau(\mathrm{L},\mathrm{M})\) does not depend on the choice of the resolution L of A/ \(\mathfrak{a}\) ; let us denote it \(\tau_{\mathfrak{a}}(\mathrm{M})\) . It then follows that the \(\tau_{\mathfrak{a}}(\mathrm{M})\) for \(\mathfrak{a} \in \mathcal{D}_{cs}\) form an inductive system of isomorphisms. Passing to the inductive limit, one obtains for each integer \(i\) , taking into account A, X, p.70, prop. 8, an isomorphism of A-modules

\[
\tau^ {i} (\mathrm{M}): \operatorname{Tor} _ {d - i} ^ {\mathrm{A}} (\mathrm{M}, \mathrm{H} _ {\mathrm{A}} ^ {d} (\mathrm{A})) \longrightarrow \mathrm{H} _ {\mathrm{A}} ^ {i} (\mathrm{M}).
\]

For M = A, \(\tau^{d}(A)\) is the canonical isomorphism from \(A \otimes_{A} H_{A}^{d}(A)\) onto \(H_{A}^{d}(A)\) .

Remarks.--- 1) Let \(f: \mathrm{M} \to \mathrm{N}\) be a homomorphism of A-modules. Using A, X, p.103, prop. 2, one proves that the following diagrams are commutative:

\[
\begin{array}{c c c} \operatorname{Tor} _ {d - i} ^ {\mathrm{A}} (\mathrm{M}, \mathrm{H} _ {\mathrm{A}} ^ {d} (\mathrm{A})) & \xrightarrow {\tau^ {i} (\mathrm{M})} & \mathrm{H} _ {\mathrm{A}} ^ {i} (\mathrm{M}) \\ \operatorname{Tor} _ {d - i} (f, 1) \Bigg \downarrow & & \Bigg \downarrow \mathrm{H} ^ {i} (f) \\ \operatorname{Tor} _ {d - i} ^ {\mathrm{A}} (\mathrm{N}, \mathrm{H} _ {\mathrm{A}} ^ {d} (\mathrm{A})) & \xrightarrow {\tau^ {i} (\mathrm{N})} & \mathrm{H} _ {\mathrm{A}} ^ {i} (\mathrm{N}) \quad . \end{array}
\]

\begin{enumerate}
\def\labelenumi{\arabic{enumi})}
\setcounter{enumi}{1}
\tightlist
\item
  Let
\end{enumerate}

\[
0 \to \mathrm{M} \to \mathrm{N} \to \mathrm{P} \to 0\tag{ℰ}
\]

an exact sequence of A-modules. Using A, X, p.104, prop. 3 and p.106, prop. 4, one proves that the following diagrams are commutative:

\[
\begin{array}{c c c} \operatorname{Tor} _ {d - i} ^ {\mathrm{A}} (\mathrm{P}, \mathrm{H} _ {\mathrm{A}} ^ {d} (\mathrm{A})) & \xrightarrow {\tau^ {i} (\mathrm{P})} & \mathrm{H} _ {\mathrm{A}} ^ {i} (\mathrm{P}) \\ \partial_ {d - i} (\mathscr {E}, \mathrm{H} _ {\mathrm{A}} ^ {d} (\mathrm{A})) \Bigg \downarrow & & \Bigg \downarrow \partial^ {i} (\mathscr {E}) \\ \operatorname{Tor} _ {d - i - 1} ^ {\mathrm{A}} (\mathrm{M}, \mathrm{H} _ {\mathrm{A}} ^ {d} (\mathrm{A})) & \xrightarrow {\tau^ {i + 1} (\mathrm{M})} & \mathrm{H} _ {\mathrm{A}} ^ {i + 1} (\mathrm{M}) \end{array} .
\]

\begin{enumerate}
\def\labelenumi{\arabic{enumi})}
\setcounter{enumi}{2}
\tightlist
\item
  Consider the isomorphism
\end{enumerate}

\[
\tau^ {d} (\mathrm{M}): \mathrm{M} \otimes_ {\mathrm{A}} \mathrm{H} _ {\mathrm{A}} ^ {d} (\mathrm{A}) \longrightarrow \mathrm{H} _ {\mathrm{A}} ^ {d} (\mathrm{M}).
\]

For \(x \in \mathrm{M}\) let us denote by \(f_x\) the map \(a \mapsto ax\) from A into M. For every \(u \in \mathrm{H}_{\mathrm{A}}^{d}(\mathrm{A})\) , one has

\[
\tau^ {d} (\mathrm{M}) (x \otimes u) = \mathrm{H} _ {\mathrm{A}} ^ {d} (f _ {x}) (u).
\]

This indeed follows from Remark 1 applied to the homomorphism \(f_{x}:A\to M\) .

\subsection{Grothendieck duality over a Macaulay ring}

We continue to assume that A is a local Macaulay ring, of dimension \(d\) . Let \(\Omega\) be a dualizing A-module. The A-module \(\mathrm{H}_{\mathrm{A}}^{d}(\Omega)\) is a Matlis module ( \(\S 8\) , no. 5, Example 6); in accordance with the notation of \(\S 8\) , we set \(\mathrm{D(M)} = \mathrm{Hom}_{\mathrm{A}}(\mathrm{M}, \mathrm{H}_{\mathrm{A}}^{d}(\Omega))\) for every A-module M.

Consider the isomorphism \(\tau^{d}(\Omega):\Omega\otimes\mathrm{H}_{\mathrm{A}}^{d}(\mathrm{A})\to\mathrm{H}_{\mathrm{A}}^{d}(\Omega)\) (no. 2). From this one deduces a homomorphism \(\omega:\mathrm{H}_{\mathrm{A}}^{d}(\mathrm{A})\to\mathrm{Hom}_{\mathrm{A}}(\Omega,\mathrm{H}_{\mathrm{A}}^{d}(\Omega))\) which associates with each element \(u\) of \(\mathrm{H}_{\mathrm{A}}^{d}(\mathrm{A})\) the homomorphism \(x\mapsto\mathrm{H}_{\mathrm{A}}^{d}(f_{x})(u)\) (Remark 3 above).

PROPOSITION 2.--- Let A be a local Macaulay ring of dimension d, and let \(\Omega\) be a dualizing A-module. The homomorphism \(\omega: H_{A}^{d}(A) \to D(\Omega)\) is bijective.

The homomorphism \(\omega\) is the limit of the inductive system of maps \((\omega_{\mathfrak{a}})_{\mathfrak{a}\in\mathcal{D}_{cs}}\) , where

\[
\omega_ {\mathfrak {a}}: \operatorname{Ext} _ {\mathrm{A}} ^ {d} (\mathrm{A} / \mathfrak {a}, \mathrm{A}) \longrightarrow \operatorname{Hom} _ {\mathrm{A}} (\Omega , \operatorname{Ext} _ {\mathrm{A}} ^ {d} (\mathrm{A} / \mathfrak {a}, \Omega))
\]

which associates with each element \(u\) of \(\operatorname{Ext}_{\mathrm{A}}^{d}(\mathrm{A}/\mathfrak{a},\mathrm{A})\) the map \(x\mapsto f_{x}\circ u\) (A, X, p.114). It therefore suffices to prove that each of the maps \(\omega_{\mathfrak{a}}\) is bijective.

Let \(\mathfrak{a}\) be an ideal of \(\mathcal{D}_{cs}\) generated by a sequence \(\mathbf{x} = (x_1, \ldots, x_d)\) completely secant for A. The Koszul complex \(\mathbf{K}^{\bullet}(\mathbf{x}, \mathrm{A})\) provides a projective resolution of A/ \(\mathfrak{a}\) ; for every A-module M, the A-module \(\mathrm{H}^d\big(\mathrm{Homgr}_\mathrm{A}(\mathbf{K}^{\bullet}(\mathbf{x}, \mathrm{A}), \mathrm{M})\big)\) is canonically identified with M/ \(\mathfrak{a}\mathrm{M}\) (A, X, p.155). From this we deduce an isomorphism (A, X, p.100)

\[
\varphi_ {\mathrm{M}}: \mathrm{M} / \mathfrak {a} \mathrm{M} \longrightarrow \operatorname{Ext} ^ {d} (\mathrm{A} / \mathfrak {a}, \mathrm{M}).
\]

Let \(x \in \Omega\) . In view of loc. cit., p.103, prop. 2, we have a commutative diagram

\[
\begin{array}{c c c} \mathrm{A/a} & \xrightarrow {\varphi_ {\mathrm{A}}} & \operatorname{Ext} _ {\mathrm{A}} ^ {d} (\mathrm{A/a}, \mathrm{A}) \\ \bar {f} _ {x} \Bigg \downarrow & & \Bigg \downarrow \operatorname{Ext} (1 _ {\mathrm{A} / a}, f _ {x}) \\ \Omega / a \Omega & \xrightarrow {\varphi_ {\Omega}} & \operatorname{Ext} _ {\mathrm{A}} ^ {d} (\mathrm{A/a}, \Omega) \end{array}
\]

where \(\bar{f}_x\) is the homomorphism deduced from \(f_x\) by passing to quotients. It follows that if for every A-module M one identifies \(\mathrm{Ext}_{\mathrm{A}}^{d}(\mathrm{A}/\mathfrak{a},\mathrm{M})\) with M/ \(\mathfrak{a}\) M by means of \(\varphi_{\mathrm{M}}\) , the homomorphism \(\omega_{\mathfrak{a}}\) is identified with the A-linear map from A/ \(\mathfrak{a}\) to \(\mathrm{Hom}_{\mathrm{A}}(\Omega,\Omega/\mathfrak{a}\Omega)\) which sends 1 to the canonical surjection, that is, to the canonical map A/ \(\mathfrak{a}\longrightarrow\mathrm{End}_{\mathrm{A}/\mathfrak{a}}(\Omega/\mathfrak{a}\Omega)\) . But since the A/ \(\mathfrak{a}\) -module \(\Omega/\mathfrak{a}\Omega\) is dualizing (\(\S\) 9, no. 2, prop. 4), this map is bijective ( \(\S\) 9, no. 4, prop. 6), which proves the proposition.

Identify the Matlis double dual D(D(Ω)) with Ω via the isomorphism \(\widehat{\alpha}_{\Omega}\) (§ 8, no. 3, th. 2, b)).

COROLLARY.--- The homomorphism D(ω): \(\widehat{\Omega} \rightarrow \mathrm{D}(\mathrm{H}_{\mathrm{A}}^{d}(\mathrm{A}))\) is an isomorphism.

Let M be an A-module, and let i be an integer. Consider the canonical homomorphisms (§ 8, no. 7)

\[
\rho_ {d - i} (\mathrm{M}, \Omega): \mathrm{Tor} _ {d - i} ^ {\mathrm{A}} (\mathrm{M}, \mathrm{D} (\Omega)) \longrightarrow \mathrm{D} \big (\mathrm{Ext} _ {\mathrm{A}} ^ {d - i} (\mathrm{M}, \Omega) \big)
\]

\[
\theta^ {d - i} (\mathrm{M}, \mathrm{H} _ {\mathrm{A}} ^ {d} (\mathrm{A})): \operatorname{Ext} _ {\mathrm{A}} ^ {d - i} \left(\mathrm{M}, \mathrm{D} \left(\mathrm{H} _ {\mathrm{A}} ^ {d} (\mathrm{A})\right)\right) \longrightarrow \mathrm{D} \left(\operatorname{Tor} _ {d - i} ^ {\mathrm{A}} \left(\mathrm{M}, \mathrm{H} _ {\mathrm{A}} ^ {d} (\mathrm{A})\right)\right).
\]

Using the isomorphisms \(\omega : \mathrm{H}_{\mathrm{A}}^{d}(\mathrm{A}) \to \mathrm{D}(\Omega)\) , \(\mathrm{D}(\omega) : \widehat{\Omega} \to \mathrm{D}(\mathrm{H}_{\mathrm{A}}^{d}(\mathrm{A}))\) (corollary 1 of Proposition 2) and \(\tau^{i}(\mathrm{M}) : \mathrm{Tor}_{d-i}^{\mathrm{A}}(\mathrm{M}, \mathrm{H}_{\mathrm{A}}^{d}(\mathrm{A})) \longrightarrow \mathrm{H}_{\mathrm{A}}^{i}(\mathrm{M})\) ( \(n^{\circ}\) 2), we deduce canonical homomorphisms of A-modules

\[
\gamma^ {i} (\mathrm{M}): \mathrm{H} _ {\mathrm{A}} ^ {i} (\mathrm{M}) \longrightarrow \mathrm{D} \left(\mathrm{Ext} _ {\mathrm{A}} ^ {d - i} (\mathrm{M}, \Omega)\right)
\]

\[
\delta^ {i} (\mathrm{M}): \operatorname{Ext} _ {\mathrm{A}} ^ {d - i} (\mathrm{M}, \widehat {\Omega}) \longrightarrow \mathrm{D} \left(\mathrm{H} _ {\mathrm{A}} ^ {i} (\mathrm{M})\right).
\]

THEOREM 1 (Grothendieck duality). Let A be a local Macaulay ring of dimension d, and let Ω be a dualizing A-module.

\begin{enumerate}
\def\labelenumi{\alph{enumi})}
\item
  The A-module \(\mathrm{H}_{\mathrm{A}}^{d}(\Omega)\) is a Matlis module; for every A-module P, let D(P) denote the Matlis dual \(\mathrm{Hom}_{\mathrm{A}}(\mathrm{P},\mathrm{H}_{\mathrm{A}}^{d}(\Omega))\) .
\item
  For every \(\mathrm{A}\)-module of finite type M and any integer \(i\) , the canonical homomorphism
\end{enumerate}

\[
\gamma^ {i} (\mathrm{M}): \mathrm{H} _ {\mathrm{A}} ^ {i} (\mathrm{M}) \longrightarrow \mathrm{D} \left(\operatorname{Ext} _ {\mathrm{A}} ^ {d - i} (\mathrm{M}, \Omega)\right)
\]

is an isomorphism of Artinian A-modules.

\begin{enumerate}
\def\labelenumi{\alph{enumi})}
\setcounter{enumi}{2}
\tightlist
\item
  For every A-module M and every integer i, the canonical homomorphism
\end{enumerate}

\[
\delta^ {i} (\mathrm{M}): \operatorname{Ext} _ {\mathrm{A}} ^ {d - i} (\mathrm{M}, \widehat {\Omega}) \longrightarrow \mathrm{D} \left(\mathrm{H} _ {\mathrm{A}} ^ {i} (\mathrm{M})\right)
\]

is an isomorphism of \(\widehat{\mathrm{A}}\) -modules.

This follows from prop. 6 of § 8, no. 7.

COROLLARY.- Let A be a local Macaulay ring, M a nonzero finitely generated A-module of dimension e. The A-module \(\mathrm{H}_{\mathrm{A}}^{e}(\mathrm{M})\) is nonzero.

By virtue of Remark 2 of No. 1, we may assume that the local ring A is complete. In this case A possesses a dualizing module Ω (§ 9, No. 3, cor. 3 of prop. 6); if \(\mathrm{H}_{\mathrm{A}}^{e}(M)\) is zero, then so is its Matlis dual \(\operatorname{Ext}_{\mathrm{A}}^{d-e}(M, \Omega)\), which contradicts Prop. 3, b) of § 9, No. 1.

Remarks.--- 1) When the A-module M is finitely generated, the \(\widehat{\mathrm{A}}\) -module \(\operatorname{Ext}_{\mathrm{A}}^{d-i}(\mathrm{M}, \widehat{\Omega})\) is identified with \(\widehat{\mathrm{A}} \otimes_{\mathrm{A}} \operatorname{Ext}_{\mathrm{A}}^{d-i}(\mathrm{M}, \Omega)\) (A, X, p.108, prop. 7, c)), and \(\delta^{i}(\mathrm{M})\) can also be obtained by composing \(\mathrm{D}(\gamma^{i}(\mathrm{M}))\) with the biduality isomorphism.

\begin{enumerate}
\def\labelenumi{\arabic{enumi})}
\setcounter{enumi}{1}
\tightlist
\item
  Let \(u: \mathrm{M} \to \mathrm{M}'\) be a homomorphism of A-modules. By Remark 1 of No. 2 and that of § 8, No. 7, the following diagrams are commutative:
\end{enumerate}

\[
\begin{array}{c c c} \mathrm{H} _ {\mathrm{A}} ^ {i} (\mathrm{M}) & \xrightarrow {\gamma^ {i} (\mathrm{M})} & \mathrm{D} (\mathrm{Ext} _ {\mathrm{A}} ^ {d - i} (\mathrm{M}, \Omega)) \\ \mathbb {1} _ {\mathrm{A}} ^ {i} (u) \Bigg \downarrow & & \Bigg \downarrow \mathrm{D} (\mathrm{Ext} (u, 1)) \\ \mathrm{H} _ {\mathrm{A}} ^ {i} (\mathrm{M} ^ {\prime}) & \xrightarrow {\gamma^ {i} (\mathrm{M} ^ {\prime})} & \mathrm{D} (\mathrm{Ext} _ {\mathrm{A}} ^ {d - i} (\mathrm{M} ^ {\prime}, \Omega)) \end{array}
\]

\[
\begin{array}{c c c} \operatorname{Ext} _ {\mathrm{A}} ^ {d - i} (\mathrm{M} ^ {\prime}, \widehat {\Omega}) & \xrightarrow {\delta^ {i} (\mathrm{M} ^ {\prime})} & \mathrm{D} (\mathrm{H} _ {\mathrm{A}} ^ {i} (\mathrm{M} ^ {\prime})) \\ \operatorname{Ext} (u, 1) \Bigg \downarrow & & \Bigg \downarrow \mathrm{D} (\mathrm{H} _ {\mathrm{A}} ^ {i} (u)) \\ \operatorname{Ext} _ {\mathrm{A}} ^ {d - i} (\mathrm{M}, \widehat {\Omega}) & \xrightarrow {\delta^ {i} (\mathrm{M})} & \mathrm{D} (\mathrm{H} _ {\mathrm{A}} ^ {i} (\mathrm{M})) \end{array} .
\]

\begin{enumerate}
\def\labelenumi{\arabic{enumi})}
\setcounter{enumi}{2}
\tightlist
\item
  Let
\end{enumerate}

\[
0 \rightarrow \mathrm{M} ^ {\prime} \rightarrow \mathrm{M} \rightarrow \mathrm{M} ^ {\prime \prime} \rightarrow 0\tag{ℰ}
\]

an exact sequence of A-modules. By Remark 2 of No. 2 and that of § 8, No. 7, the following diagrams are commutative:

\[
\begin{array}{c c c} \mathrm{H} _ {\mathrm{A}} ^ {i - 1} (\mathrm{M} ^ {\prime \prime}) & \xrightarrow {\gamma^ {i - 1} (\mathrm{M} ^ {\prime \prime})} & \mathrm{D} (\mathrm{Ext} _ {\mathrm{A}} ^ {d - i + 1} (\mathrm{M} ^ {\prime \prime}, \Omega)) \\ \partial^ {i - 1} (\mathscr {E}) \Bigg \downarrow & & \Bigg \downarrow (- 1) ^ {d - i + 1} \mathrm{D} (\delta^ {d - i} (\mathscr {E}, \Omega)) \\ \mathrm{H} _ {\mathrm{A}} ^ {i} (\mathrm{M} ^ {\prime}) & \xrightarrow {\gamma^ {i} (\mathrm{M} ^ {\prime})} & \mathrm{D} (\mathrm{Ext} _ {\mathrm{A}} ^ {d - i} (\mathrm{M} ^ {\prime}, \Omega)) \end{array}
\]

\[
\begin{array}{c c c} \operatorname{Ext} _ {\mathrm{A}} ^ {d - i} (\mathrm{M} ^ {\prime}, \widehat {\Omega}) & \xrightarrow {\delta^ {i} (\mathrm{M} ^ {\prime})} & \mathrm{D} (\mathrm{H} _ {\mathrm{A}} ^ {i} (\mathrm{M} ^ {\prime})) \\ \delta^ {d - i} (\mathscr {E}, \widehat {\Omega}) \Bigg \downarrow & & \Bigg \downarrow (- 1) ^ {d - i + 1} \mathrm{D} (\partial^ {i - 1} (\mathscr {E})) \\ \operatorname{Ext} _ {\mathrm{A}} ^ {d - i + 1} (\mathrm{M} ^ {\prime \prime}, \widehat {\Omega}) & \xrightarrow {\delta^ {i - 1} (\mathrm{M} ^ {\prime \prime})} & \mathrm{D} (\mathrm{H} _ {\mathrm{A}} ^ {i - 1} (\mathrm{M} ^ {\prime \prime})) \quad . \end{array}
\]

Example.--- Let A be a Noetherian local integral domain of dimension 1; denote by K its field of fractions. Let Ω be a dualizing A-module, and let M be a finitely generated A-module. The A-modules \(\mathrm{H}_{\mathrm{A}}^{0}(M)\) and \(\mathrm{H}_{\mathrm{A}}^{1}(M)\) are canonically identified with T(M) and \((K/A) \otimes_{A} M\) (no. 1, example 5). With these identifications, the duality isomorphisms

\[
\gamma^ {0} (\mathrm{M}): \mathrm{T} (\mathrm{M}) \longrightarrow \mathrm{D} \left(\operatorname{Ext} _ {\mathrm{A}} ^ {1} (\mathrm{M}, \Omega)\right), \quad \gamma^ {1} (\mathrm{M}): (\mathrm{K} / \mathrm{A}) \otimes_ {\mathrm{A}} \mathrm{M} \longrightarrow \mathrm{D} \left(\operatorname{Hom} _ {\mathrm{A}} (\mathrm{M}, \Omega)\right)
\]

(th. 1) are none other than the isomorphisms defined in § 9, no. 6.

\specialchapter{Exercises}

\exercisesection{Depth}

\begin{enumerate}
\def\labelenumi{\arabic{enumi})}
\item
  Let A be a ring, J an ideal of A, and M an A-module. Prove that in order to have \(\operatorname{prof}_{\mathrm{A}}(J;M)\geqslant2\) it is necessary and sufficient that the map \(M\to\operatorname{Hom}_{\mathrm{A}}(J,M)\) deduced from the canonical injection of J into A be an isomorphism.
\item
  Let A be a valuation ring of height 1, not Noetherian, m its maximal ideal, n a principal ideal of A, distinct from (0) and from A. Prove that we have \(\operatorname{prof}_{\mathrm{A}}(\mathfrak{n};\mathrm{A})=1\) and \(\operatorname{prof}_{\mathrm{A}}(\mathfrak{m};\mathrm{A})\geqslant2\) even though one has \(\mathrm{V}(\mathfrak{m})=\mathrm{V}(\mathfrak{n})\) (Observe that \(\operatorname{Ext}_{\mathrm{A}}^{1}(\mathrm{A}/\mathfrak{m},\mathrm{A}/\mathfrak{n})\) is isomorphic to \(\operatorname{Hom}_{\mathrm{A}}(\mathrm{A}/\mathfrak{m},\mathrm{A})\) .)
\item
  Let A be the local ring \(k[[X,Y]]\) and let M be the direct sum of the A-modules \(\mathrm{A}/\mathfrak{a}\) , where \(\mathfrak{a}\) runs through the set of nonzero principal ideals of A. Show that one has \(\mathrm{prof}_{\mathrm{A}}(\mathrm{M}) = 1\) but that there exists no M-regular element in \(\mathfrak{m}_{\mathrm{A}}\) .
\item
  Let A be a ring, J a finitely generated ideal of A, M an A-module, P a flat A-module. Prove that one has \(\operatorname{prof}_{\mathrm{A}}(\mathrm{J};\mathrm{P}\otimes_{\mathrm{A}}\mathrm{M})\geqslant\operatorname{prof}_{\mathrm{A}}(\mathrm{J};\mathrm{M})\) and that equality holds if P is faithfully flat.
\item
  Let \(\mathrm{A}\) be a Noetherian ring, M a nonzero finitely generated A-module, J an ideal of \(\mathrm{A}\) , \(\mathbf{x} = (x_1, \ldots, x_n)\) a generating system of J. Prove that we have \(\mathrm{H}^i(\mathbf{x}, \mathrm{M}) \neq 0\) for \(\operatorname{prof}_{\mathrm{A}}(\mathrm{J}; \mathrm{M}) \leqslant i \leqslant n\) (reduce to the local case, and use A, X, p.157, cor. 2).
\item
  Let A be a Noetherian local ring, and let P be a complex of flat A-modules of finite length \(\ell\) , such that \(\operatorname{Supp}(\mathrm{H}(\mathrm{P})) = \{\mathfrak{m}_{\mathrm{A}}\}\) . Prove that every A-module M such that \(\kappa_{\mathrm{A}} \otimes_{\mathrm{A}} \mathrm{M} \neq 0\) is of depth \(\leqslant \ell\) (applying prop. 3 to the complex \(\mathrm{P} \otimes_{\mathrm{A}} \mathrm{M}\) by observing that this complex is not exact and using exerc. 4).
\item
  Let A be a local ring, M a finitely generated A-module, and F a closed subset of \(\operatorname{Spec}(A)\) . Prove the inequality \(\operatorname{prof}_{\mathbf{F}}(\mathbf{M}) \geqslant \operatorname{prof}(\mathbf{M}) - \dim(\mathbf{F})\) .
\item
  Let A be a Noetherian ring, and let M be a finitely generated A-module. We say that M satisfies the property \((S_{k})\) if one has \(\operatorname{prof}_{A_{\mathfrak{p}}}(M_{\mathfrak{p}}) \geqslant \inf(k, \dim_{A_{\mathfrak{p}}}(M_{\mathfrak{p}}))\) for every prime ideal p of A.
\end{enumerate}

\begin{enumerate}
\def\labelenumi{\alph{enumi})}
\item
  Every module satisfies \((\mathrm{S}_{k})\) for \(k \leqslant 0\) ; to say that a module satisfies \((\mathrm{S}_{1})\) means that it has no embedded associated prime ideals.
\item
  Let \(0 \to M' \to M \to M'' \to 0\) be a short exact sequence of finitely generated A-modules, and \(k\) an integer. We assume that \(M\) satisfies \((S_k)\) , that \(M''\) satisfies \((S_{k-1})\) , and if \(k \geqslant 2\) that we have \(\operatorname{Ass}(M'') \subset \operatorname{Ass}(M)\) . Prove that \(M'\) satisfies \((S_k)\) (using cor. 2 of no. 7).
\item
  c) Let \((L,p)\) be a resolution of M by finitely generated free A-modules. If A satisfies the property \((S_{k})\) , the A-module \(B_{i}(L)\) for \(i \geqslant 0\) satisfies \((S_{h})\) with \(h = \inf(k,i+1)\) . d) If A satisfies \((S_{2})\) , the dual of every finitely generated A-module satisfies \((S_{2})\) (apply c). e) Let J be an ideal of A, generated by an M-regular sequence \((x_{1},\ldots,x_{r})\) . If M satisfies \((S_{k})\) , the A-module M/JM satisfies \((S_{k-r})\) .
\end{enumerate}

\(f)\) Let B be a Noetherian A-algebra, N a finitely generated B-module which is a faithfully flat A-module, and \(k\) an integer. If the B-module \(\mathrm{M} \otimes_{\mathrm{A}} \mathrm{N}\) satisfies property \((\mathrm{S}_k)\) , then so does the A-module M. If M satisfies \((\mathrm{S}_k)\) and if the \((\kappa(\mathfrak{p}) \otimes_{\mathrm{A}} \mathrm{B})\) -module \(\kappa(\mathfrak{p}) \otimes_{\mathrm{A}} \mathrm{N}\) satisfies \((\mathrm{S}_k)\) for every \(\mathfrak{p} \in \operatorname{Supp}(\mathrm{M})\) , the B-module \(\mathrm{M} \otimes_{\mathrm{A}} \mathrm{N}\) satisfies \((\mathrm{S}_k)\) .

\begin{enumerate}
\def\labelenumi{\arabic{enumi})}
\setcounter{enumi}{8}
\item
  Give an example of a local ring A and a finitely generated A-module M such that \(\mathfrak{m}_{\mathrm{A}}M \neq M\) and \(\operatorname{prof}_{A}(M) = +\infty\) (take for A the localization of a polynomial ring in an infinite family of indeterminates).
\item
  Let A be a Noetherian ring, M a finitely generated A-module, J an ideal of A, \((x_{1},\ldots,x_{r})\) an M-regular sequence of elements of J, with \(r < \operatorname{prof}_{\mathrm{A}}(\mathrm{J};\mathrm{M})\) . Let \(q_{1}\ldots,q_{n}\) be ideals of A not containing J, all prime except possibly two, and \(J_{0}\) a subset of J, stable under addition and multiplication and generating the ideal J. Prove that there exists an element x of \(J_{0}\) not belonging to any of the \(q_{i}\) such that the sequence \((x_{1},\ldots,x_{r},x)\) is M-regular (use Cor. 2 of Prop. 2 of II, § 1, No. 1).
\item
  Let \(\mathrm{A}\) be a ring, M an \(\mathrm{A}\)-module, and \(x_{1},\ldots ,x_{r}\) elements of \(\mathrm{A}\) .
\end{enumerate}

\begin{enumerate}
\def\labelenumi{\alph{enumi})}
\item
  For every integer p such that \(1 \leqslant p \leqslant r\) , we set \(M_{p} = M/(x_{2}M + \ldots + x_{p}M)\) . Prove that if the sequence \((x_{1}, \ldots, x_{r})\) is M-regular, the sequence \((x_{1}, x_{p+1}, x_{p+2}, \ldots, x_{r})\) is \(M_{p}\) -regular (argue by induction on p).
\item
  For the sequences \((x_{1}, x_{3})\) and \((x_{2}, x_{3})\) to be M-regular, it is necessary and sufficient that the sequence \((x_{1}x_{2}, x_{3})\) be M-regular.
\item
  Let \(x_p' \in \mathrm{A}\) ( \(1 \leqslant p \leqslant r\) ). For the sequences \((x_1, \ldots, x_p, \ldots, x_r)\) and \((x_1, \ldots, x_p', \ldots, x_r)\) to be M-regular, it is necessary and sufficient that the sequence \((x_1, \ldots, x_p x_p', \ldots, x_r)\) be M-regular (use a) and b)).
\end{enumerate}

\(d\) Let \(n_1,\ldots ,n_r\) be integers \(\geqslant 1\) . For the sequence \((x_{1}^{n_{1}},\dots,x_{r}^{n_{r}})\) to be M-regular, it is necessary and sufficient that the sequence \((x_{1},\ldots ,x_{r})\) be M-regular.

\begin{enumerate}
\def\labelenumi{\arabic{enumi})}
\setcounter{enumi}{11}
\tightlist
\item
  Let A be a ring, M an A-module.
\end{enumerate}

\begin{enumerate}
\def\labelenumi{\alph{enumi})}
\item
  Let \(P \in A[X]\) . Prove that if the kernel of the homothety of ratio P in M{[}X{]} is not zero, it contains a nonzero element of M (let Q be an element of M{[}X{]} of minimal degree such that PQ = 0; if \(P = a_{0}X^{p} + \ldots + a_{p}\) , \(Q = m_{0}X^{q} + \ldots + m_{q}\) , observe that \(a_{0}Q = 0\) , then by induction that \(a_{i}Q = 0\) for all i).
\item
  Let J be a finitely generated ideal of A. Prove that the relation \(\operatorname{prof}_{\mathrm{A}}(\mathrm{J};\mathrm{M}) > 0\) is equivalent to the existence of an element which is M{[}X{]}-regular in JA{[}X{]}.
\item
  Let I and J be finitely generated ideals of A. Prove the equality \(\mathrm{prof}_{\mathrm{A}}(\mathrm{IJ};\mathrm{M}) = \min (\mathrm{prof}_{\mathrm{A}}(\mathrm{I};\mathrm{M}),\mathrm{prof}_{\mathrm{A}}(\mathrm{J};\mathrm{M}))\) (reason by induction on \(\mathrm{prof}_{\mathrm{A}}(\mathrm{IJ};\mathrm{M})\) , using \(b)\) ).
\end{enumerate}

Let A be a ring, J an ideal of A, and M an A-module. For every integer \(n \geqslant 0\) , we denote by \(\operatorname{prof}_{n}(\mathbf{J}; \mathbf{M})\) the upper bound (in \(\overline{\mathbf{N}}\) ) of the lengths of the sequences \(\mathrm{M}[\mathrm{X}_{1}, \ldots, \mathrm{X}_{n}]\) -regular elements of \(\mathrm{JA}[\mathrm{X}_{1}, \ldots, \mathrm{X}_{n}]\) .

\begin{enumerate}
\def\labelenumi{\alph{enumi})}
\item
  Show that the sequence \(\left(\operatorname{prof}_{n}(\mathrm{J};\mathrm{M})\right)_{n\geqslant0}\) is increasing; we denote by \(\operatorname{prof}_{\infty}(\mathrm{J};\mathrm{M})\) its limit (in \(\overline{\mathbf{N}}\) ). Prove that we have \(\operatorname{prof}_{\infty}(\mathrm{J};\mathrm{M})\leqslant\operatorname{prof}_{\mathrm{A}}(\mathrm{J};\mathrm{M})\) .
\item
  When the ideal J is finitely generated, prove the equality \(\operatorname{prof}_{\infty}(\mathrm{J};\mathrm{M})=\operatorname{prof}_{\mathrm{A}}(\mathrm{J};\mathrm{M})\) (reason as in the proof of Th. 2 of No. 4, using Exerc. 12.)
\item
  In the general case, show that \(\operatorname{prof}_{\infty}(J;M)\) is the supremum of the numbers \(\operatorname{prof}_{\mathrm{A}}(J';M)\) , where \(J'\) ranges over the set of finitely generated ideals contained in J.
\item
  Let \(J'\) be an ideal of A such that \(V(J') = V(J)\) . Prove that we have \(\text{prof}_{\infty}(J'; M) = \text{prof}_{\infty}(J; M)\) .
\item
  Let A be a valuation ring of height 1, not Noetherian, and let m be its maximal ideal. Prove that we have \(\operatorname{prof}_{\infty}(\mathfrak{m};\mathrm{A})=1\) and \(\operatorname{prof}_{\mathrm{A}}(\mathfrak{m};\mathrm{A})\geqslant2\) (cf. exerc. 2).
\end{enumerate}

\begin{enumerate}
\def\labelenumi{\arabic{enumi})}
\setcounter{enumi}{13}
\tightlist
\item
  Let A be a complete Noetherian local ring.
\end{enumerate}

\begin{enumerate}
\def\labelenumi{\alph{enumi})}
\item
  Prove that an ideal of A contained in the union of a sequence of prime ideals of A is contained in one of them (observe that the topological space A is a Baire space, cf. TG, IX, p. 55, th. 1).
\item
  Let M be an A-module admitting a countable generating family. Show that the conclusions of Th. 2 still hold (observe that the set Ass(M) is countable, and deduce from a) that if \(\mathrm{prof}_{\mathrm{A}}(\mathrm{J};\mathrm{M}) > 0\) , there exists an element \(x\) of J such that the homothety \(x_{\mathrm{M}}\) is injective).
\end{enumerate}

\begin{enumerate}
\def\labelenumi{\arabic{enumi})}
\setcounter{enumi}{14}
\tightlist
\item
  Let A be a Noetherian local ring, and let M be a finitely generated A-module. Prove the inequalities
\end{enumerate}

\[
\operatorname{prof} (A) \leqslant \operatorname{grade} (M) + \dim (M) \leqslant \dim (A).
\]

(Prove the first inequality by induction on the integer \(g = \text{grade}(M)\) using cor. 2 of prop. 13, no. 7 in the case \(g = 0\) and then by considering an A-regular element of Ann(M). For the second, observe that one has grade(M) ≤ prof(A \(_{p}\) ) for every p ∈ Supp(M).)

\begin{enumerate}
\def\labelenumi{\arabic{enumi})}
\setcounter{enumi}{15}
\tightlist
\item
  Let A be a Noetherian local ring. For every finitely generated A-module M, we denote by \(t_{\mathrm{A}}(\mathrm{M})\) , or simply \(t(\mathrm{M})\) , the dimension of the \(\kappa_{A}\) -vector space \(\operatorname{Ext}_{\mathrm{A}}^{p}(\kappa_{\mathrm{A}},\mathrm{M})\) , with \(p=\operatorname{prof}(\mathrm{M})\) .
\end{enumerate}

\begin{enumerate}
\def\labelenumi{\alph{enumi})}
\item
  Let \((x_{1},\ldots,x_{r})\) be a sequence of elements of \(m_{A}\) completely secant for M, generating an ideal J. We have \(t(\mathrm{M})=t(\mathrm{M}/\mathrm{JM})\) .
\item
  Let \(\rho : A \to B\) be a local homomorphism of Noetherian local rings; let M be a finitely generated A-module, and N a finitely generated B-module flat over A. Prove the formula \(t_{B}(M \otimes_{A} N) = t_{A}(M)t_{B}(\kappa_{A} \otimes_{A} N)\) (reduce using \(a\) to the case where \(\mathrm{prof}_{\mathrm{B}}(M \otimes_{A} N) = \mathrm{prof}_{\mathrm{A}}(M) = \mathrm{prof}_{\mathrm{B}}(\kappa_{\mathrm{A}} \otimes_{\mathrm{A}} N) = 0\) ).
\end{enumerate}

\begin{enumerate}
\def\labelenumi{\arabic{enumi})}
\setcounter{enumi}{16}
\tightlist
\item
  Let A be a Noetherian ring, M a finitely generated A-module. Show that M is reflexive if and only if the following two conditions are satisfied:
\end{enumerate}

\begin{enumerate}
\def\labelenumi{(\roman{enumi})}
\item
  for every \(\mathfrak{p} \in \operatorname{Spec}(A)\) such that \(\operatorname{prof}(A_{\mathfrak{p}}) \leqslant 1\) , the \(A_{\mathfrak{p}}\) -module \(M_{\mathfrak{p}}\) is reflexive;
\item
  for every \(\mathfrak{p} \in \operatorname{Spec}(\mathrm{A})\) such that \(\operatorname{prof}(A_{\mathfrak{p}}) \geqslant 2\) , one has \(\operatorname{prof}_{\mathrm{A}_{\mathfrak{p}}} (M_{\mathfrak{p}}) \geqslant 2\) .
\end{enumerate}

(Under hypotheses (i) and (ii), prove that the kernel, then the cokernel, of the canonical homomorphism \(M \to M^{**}\) have no associated prime ideals.)

\begin{enumerate}
\def\labelenumi{\arabic{enumi})}
\setcounter{enumi}{17}
\tightlist
\item
  Let A be a Noetherian ring, J an ideal of A, M and N finitely generated A-modules, with \(\operatorname{prof}_{\mathrm{A}}(J;N)\geqslant2\) .
\end{enumerate}

\begin{enumerate}
\def\labelenumi{\alph{enumi})}
\item
  Prove the inequality \(\mathrm{prof}_{\mathrm{A}}(\mathrm{J};\mathrm{Hom}_{\mathrm{A}}(\mathrm{M},\mathrm{N}))\geqslant 2\) (consider an N-regular sequence \((x,y)\) in J).
\item
  Assume \(\mathrm{prof}_{\mathrm{A}}(\mathrm{J};\mathrm{Hom}_{\mathrm{A}}(\mathrm{M},\mathrm{N}))\geqslant 3\) . Prove that we have \(\mathrm{prof}_{\mathrm{A}}(\mathrm{J};\mathrm{Ext}_{\mathrm{A}}^{1}(\mathrm{M},\mathrm{N}))\geqslant 1\) (consider an exact sequence of extension modules associated with an N-regular element of J).
\end{enumerate}

\exercisesection{Macaulay modules and rings}

\begin{enumerate}
\def\labelenumi{\arabic{enumi})}
\item
  Let A be a Noetherian local ring, and let \(0 \to M' \to M \to M'' \to 0\) be an exact sequence of Macaulay A-modules. Show that we have \(\dim(M') = \dim(M)\) and that \(\dim(M'')\) is equal to \(\dim(M)\) or to \(\dim(M) - 1\) .
\item
  Let A be a Macaulay local ring, and \(\mathfrak{p} \in \operatorname{Spec}(A)\) . Show that we have \(\dim(\widehat{A}/\mathfrak{q}) = \dim(A/\mathfrak{p})\) for every \(\mathfrak{q} \in \operatorname{Ass}(\widehat{A}/\mathfrak{p}\widehat{A})\) . In particular, the ring \(\widehat{\mathrm{A}}/\mathfrak{p}\widehat{A}\) has no embedded prime ideal.
\item
  \begin{enumerate}
  \def\labelenumii{\alph{enumii})}
  \tightlist
  \item
    Let R be a Noetherian, integral, integrally closed Q-algebra, A a ring, and \(\rho: \mathrm{R} \to \mathrm{A}\) an injective homomorphism making A into a finite R-algebra. Prove that the sub-R-module \(\rho(\mathrm{R})\) is a direct summand of A (reduce to cor. 3 of prop. 8).
  \end{enumerate}
\end{enumerate}

\begin{enumerate}
\def\labelenumi{\alph{enumi})}
\setcounter{enumi}{1}
\tightlist
\item
  Let A be a Noetherian local Q-algebra. Prove that A satisfies the following property (the “monomial condition”):
\end{enumerate}

\begin{enumerate}
\def\labelenumi{(\Roman{enumi})}
\setcounter{enumi}{899}
\tightlist
\item
  for every maximal secant sequence \((x_{1},\ldots,x_{d})\) of elements of \(m_{A}\) and every integer p, the element \((x_{1}\ldots x_{d})^{p}\) does not belong to the ideal \((x_{1}^{p+1},\ldots,x_{d}^{p+1})\) . (Reduce to the case where A is complete, hence admits a coefficient field K, and apply a) to a homomorphism \(\rho:K[[X_{1},\ldots,X_{d}]]\to A\) such that \(\rho(X_{i})=x_{i}\) .)
\end{enumerate}

\begin{enumerate}
\def\labelenumi{\arabic{enumi})}
\setcounter{enumi}{3}
\tightlist
\item
  Let A be a local ring, p a prime ideal of A. Let B denote the ring of fractions \(A[Z][S^{-1}]\) , where S is the set of elements \(Z + a\) of A{[}Z{]} with \(a \in \mathfrak{m}_{A} - \mathfrak{p}\) . Let z denote the element Z/1 of B.
\end{enumerate}

\begin{enumerate}
\def\labelenumi{\alph{enumi})}
\item
  Prove that \(\mathrm{B} / (z)\) is isomorphic to \(\mathbf{A}_{\mathfrak{p}}\) and that \(\mathrm{B} / (z - 1)\) is isomorphic to \(\mathbf{A}\) .
\item
  Assume henceforth that the ring A is integral and Noetherian, and that p is generated by a nonzero element p of A. Let C be the ring B{[}T{]}/(z(pT-1)); if A is a Macaulay ring, so is C. Prove that Spec C has two irreducible components \(X_{1}\) and \(X_{2}\) which meet at a closed point x, so that \(\text{codim}(\{x\}, X_{1}) = \text{codim}(\{x\}, X_{2}) = 1\) , \(\text{dim}(X_{1}) = 2\) and \(\text{dim}(X_{2}) \geqslant \text{dim}(A_{\mathfrak{p}})\) . For every integer n, give an example with \(\text{dim}(A_{\mathfrak{p}}) = n\) .
\end{enumerate}

¶ 5) Let A be a Noetherian local ring, M a finitely generated A-module, \(\mathbf{x} = (x_{1},\ldots ,x_{n})\) a sequence of elements of \(\mathfrak{m}_{\mathrm{A}}\) , generating an ideal \(\mathfrak{x}\); assume that the A-module M/xM has finite length (which implies \(n\geqslant \dim (\mathrm{M})\) ). Let \(\mathrm{gr(A)}\) and \(\mathrm{gr(M)}\) denote the associated graded rings of A and M for the x-adic topology, \(\xi_{i}\) ( \(1\leqslant i\leqslant n\) ) the image of \(x_{i}\) in \(\mathfrak{x} / \mathfrak{x}^{2}\) , \(\xi\) the sequence \((\xi_1,\dots ,\xi_n)\) .

\begin{enumerate}
\def\labelenumi{\alph{enumi})}
\item
  For \(p \in \mathbf{Z}\) , we set \(\mathrm{F}^p\mathbf{K}_i(\mathbf{x},\mathrm{M}) = \mathbf{K}_i(\mathbf{x},\mathfrak{x}^{p - i}\mathrm{M})\) (with the convention \(\mathfrak{x}^s = A\) for \(s \leqslant 0\) ). Prove that \(\mathrm{F}^p\mathbf{K}_{\bullet}(\mathbf{x},\mathrm{M})\) is a subcomplex of \(\mathbf{K}_{\bullet}(\mathbf{x},\mathrm{M})\) , and that the complex \(\bigoplus_{n}\mathrm{F}^p\mathbf{K}_{\bullet}(\mathbf{x},\mathrm{M}) / \mathrm{F}^{p + 1}\mathbf{K}_{\bullet}(\mathbf{x},\mathrm{M})\) is identified with \(\mathbf{K}_{\bullet}^{\mathrm{gr}(A)}(\boldsymbol {\xi},\mathrm{gr}(M))\) .
\item
  For every complex of A-modules C such that H(C) is of finite length, we set \(\chi(C) = \sum_{i}(-1)^{i}\lg(\mathrm{H}_{i}(C))\) . Prove that there exists an integer \(p_0\) such that the complex \(\mathrm{F}^{p}\mathbf{K}_{\bullet}(\mathbf{x},\mathrm{M}) / \mathrm{F}^{p + 1}\mathbf{K}_{\bullet}(\mathbf{x},\mathrm{M})\) has zero homology for \(p\geqslant p_0\) ; deduce from this the equalities
\end{enumerate}

\[
\chi \left(\mathbf {K} _ {\bullet} ^ {\operatorname{gr} (\mathrm{A})} (\boldsymbol {\xi}, \operatorname{gr} (\mathrm{M}))\right) = \chi \left(\mathbf {K} _ {\bullet} (\mathrm{x}, \mathrm{M}) / \mathrm{F} ^ {p} \mathbf {K} _ {\bullet} (\mathrm{x}, \mathrm{M})\right) = \sum (- 1) ^ {i} \binom {n} {i} \lg (\mathrm{M} / x ^ {p - i} \mathrm{M})
\]

for \(p \geqslant p_0\) .

\begin{enumerate}
\def\labelenumi{\alph{enumi})}
\setcounter{enumi}{2}
\item
  Prove that the above expression is zero if \(n > \dim(\mathrm{M})\) , and equal to \(e_{\mathfrak{x}}(\mathrm{M})\) if \(n = \dim(\mathrm{M})\) .
\item
  Let \(h \geqslant 0\) , and \(p \geqslant p_{0}\); prove that the canonical injection of \(\mathrm{F}^{p}\mathbf{K}_{\bullet}(\mathbf{x},\mathrm{M})\) into \(\mathrm{F}^{p+h}\mathbf{K}_{\bullet}(\mathbf{x},\mathrm{M})\) is a homology isomorphism. Deduce from this that \(\mathrm{H}_{i}(\mathrm{F}^{p}\mathbf{K}_{\bullet}(\mathbf{x},\mathrm{M}))\) is discrete for the x-adic topology. Prove moreover that the filtration of this module defined by the images of the maps \(\mathrm{H}_{i}(\mathrm{F}^{p+h}\mathbf{K}_{\bullet}(\mathbf{x},\mathrm{M})) \to \mathrm{H}_{i}(\mathrm{F}^{p}\mathbf{K}_{\bullet}(\mathbf{x},\mathrm{M}))\) is x-good, and deduce from this that \(\mathrm{F}^{p}\mathbf{K}_{\bullet}(\mathbf{x},\mathrm{M})\) has zero homology. Conclude that \(\chi(\mathbf{K}_{\bullet}(\mathbf{x},\mathrm{M}))\) is zero if \(n > \dim(\mathrm{M})\) , and equal to \(e_{\mathfrak{x}}(\mathrm{M})\) if \(n = \dim(\mathrm{M})\) .
\end{enumerate}

\begin{enumerate}
\def\labelenumi{\arabic{enumi})}
\setcounter{enumi}{5}
\tightlist
\item
  We keep the hypotheses of the preceding exercise; for every finitely generated A-module N such that N/xN is of finite length, we denote by \(\chi(\mathbf{x},\mathbf{N})\) the integer \(\chi(\mathbf{K}_{\bullet}(\mathbf{x},\mathbf{N}))\) .
\end{enumerate}

\begin{enumerate}
\def\labelenumi{\alph{enumi})}
\item
  For every submodule \(M'\) of \(M\) , one has \(\chi(\mathbf{x}, M) = \chi(\mathbf{x}, M') + \chi(\mathbf{x}, M/M')\) .
\item
  If \(x_{1}\mathrm{M} = 0\), show that \(\chi (\mathbf{x},\mathrm{M})\) is zero.
\item
  Denote by \(\mathbf{x}'\) the sequence \((x_2,\ldots ,x_n)\); prove the equality
\end{enumerate}

\[
\chi (\mathbf {x}, \mathrm{M}) = \chi (\mathbf {x} ^ {\prime}, \mathrm{M} / x _ {1} \mathrm{M}) - \chi (\mathbf {x} ^ {\prime}, \operatorname{Ker} (x _ {1}) _ {\mathrm{M}})
\]

(use exerc. 8 of A, X, p. 207).

\begin{enumerate}
\def\labelenumi{\alph{enumi})}
\setcounter{enumi}{3}
\tightlist
\item
  For \(1 \leqslant i \leqslant n\) , we denote \(\mathrm{K}_i\) the kernel of the homothety with ratio \(x_i\) in \(\mathrm{M}/(x_1\mathrm{M} + \ldots + x_{i-1}\mathrm{M})\) . Prove the equality
\end{enumerate}

\[
\lg (\mathrm{M} / \mathfrak {x} \mathrm{M}) - \chi (\mathbf {x}, \mathrm{M}) = \sum_ {i = 1} ^ {n} \chi ((x _ {i + 1}, \dots , x _ {n}), \mathrm{K} _ {i}).
\]

\begin{enumerate}
\def\labelenumi{\alph{enumi})}
\setcounter{enumi}{4}
\tightlist
\item
  Let \(p_1, \ldots, p_n\) be integers \(> 1\) , \(\mathbf{x}^{\mathbf{p}}\) the sequence \((x_{1}^{p_1}, \ldots, x_n^{p_n})\) . Prove the equality \(\chi(\mathbf{x}^{\mathbf{p}}, \mathrm{M}) = p_1 \ldots p_n \chi(\mathbf{x}, \mathrm{M})\) (first treat the case \(n = 1\) and then reduce to this case using cor. 2 of A, X, p. 157).
\end{enumerate}

¶ 7) Let A be a Noetherian local ring, M a finitely generated A-module. One says that a sequence \((x_1, \ldots, x_n)\) is weakly regular for M if for \(i = 1, \ldots, n\) the kernel of the homothety with ratio \(x_{i+1}\) in \(\mathrm{M}/(x_1\mathrm{M} + \ldots + x_i\mathrm{M})\) is annihilated by \(\mathfrak{m}_A\) .

\begin{enumerate}
\def\labelenumi{\alph{enumi})}
\item
  If \(n \leqslant \dim(M)\), prove that a sequence \((x_{1}, \ldots, x_{n})\) weakly regular for M is secant for M (reason by induction on n). Give an example of a weakly regular sequence that is not secant (one may take \(\mathrm{A} = M = k[X]/(X^{2})\) (where k is a field).
\item
  We say that M is a Buchsbaum module if every secant sequence for M is weakly regular for M. A Macaulay module is a Buchsbaum module; if M is a Buchsbaum module, the \(A_{\mathfrak{p}}\) -module \(M_{\mathfrak{p}}\) is Macaulay for every prime ideal \(p \neq \mathfrak{m}_{\mathrm{A}}\). c) Let M be a Buchsbaum module, \((x_{1}, \ldots, x_{n})\) a secant sequence for M; prove that \(\mathrm{M}/(x_{1}\mathrm{M} + \ldots + x_{n}\mathrm{M})\) is a Buchsbaum module.
\end{enumerate}

\(d\) ) For M to be a Buchsbaum module, it is necessary and sufficient that for every maximal secant sequence \((x_{1},\ldots ,x_{n})\) for M, one has in \(\mathrm{M}_n = \mathrm{M} / (x_1\mathrm{M} + \dots +x_{n - 1}\mathrm{M})\) the equality \(\mathrm{Ker}(x_n)_{\mathrm{M}_n} = \mathrm{Ker}(x_n^2)_{\mathrm{M}_n}\) (first treat the case \(n = 1\) , observing that \(\mathrm{Ker}(x_1)_{\mathrm{M}}\) is then of finite length and reasoning by induction on this length. In the general case, complete every secant sequence \((x_{1},\ldots ,x_{i})\) to a maximal secant sequence \((x_{1},\ldots ,x_{n})\) , and consider the sequences \((x_{1},\ldots ,x_{i - 1},x_{i + 1}^{k},\ldots ,x_{n}^{k},x_{i})\) for \(k\geqslant 1\) ).

¶ 8) Let A be a Noetherian local ring, M a finitely generated A-module. We propose to prove the equivalence of the following two conditions:

\begin{enumerate}
\def\labelenumi{(\roman{enumi})}
\item
  M is a Buchsbaum module;
\item
  there exists a positive integer \(i(M)\) such that one has \(\lg(M/qM)-e_{\mathfrak{q}}(M)=i(M)\) for every ideal \(\mathfrak{q}\subset\mathfrak{m}_{A}\) such that \(M/qM\) is of finite length.
\end{enumerate}

\begin{enumerate}
\def\labelenumi{\alph{enumi})}
\item
  Under hypothesis (i), let x be a secant sequence for M, and let x denote the ideal it generates. With the notations of exerc. 6, d), prove the equality \(\lg(\mathrm{M}/x\mathrm{M}) - e_{\mathfrak{x}}(\mathrm{M}) = \lg(\mathrm{K}_{n})\) . Let \(\mathbf{x}'\) be another secant sequence for M, and let \(\mathrm{K}_{n}^{\prime}\) be the corresponding module; prove the equality \(\lg(\mathrm{K}_{n}) = \lg(\mathrm{K}_{n}^{\prime})\) (reason by induction on \(n\) , considering an element \(t\) of \(\mathfrak{m}_{\mathrm{A}}\) such that the sequences \((x_{1},\ldots,x_{n-1},t)\) and \((x_{1}',\ldots,x_{n-1}',t)\) are secant for M).
\item
  We assume hypothesis (ii) is satisfied. Let \(\mathbf{x} = (x_{1}, \ldots, x_{n})\) be a maximal secant sequence for M, and let \(M_{n} = \mathrm{M}/(x_{1}\mathrm{M} + \ldots + x_{n-1}\mathrm{M})\) ; by applying the formula of exerc. 6, d) to the sequences x and \((x_{1}, \ldots, x_{n-1}, x_{n}^{2})\) (taking account of exerc. 6, e)), prove that one has \(\operatorname{Ker}(x_{n})_{\mathrm{M}_{n}} = \operatorname{Ker}(x_{n}^{2})_{\mathrm{M}_{n}}\) . Conclude with exerc. 7, d)).
\item
  For M to be a Buchsbaum module, it is necessary and sufficient that it be so for \(\widehat{\mathbf{M}}\) ; in this case we have \(i(\widehat{\mathbf{M}}) = i(\mathbf{M})\) .
\item
  Let M be a Buchsbaum module, and \((x_{1},\ldots,x_{p})\) be a sequence of elements of \(\mathfrak{m}_{A}\) secant for M, with \(p<\dim(M)\) ; prove that we have \(i(\mathrm{M}/(x_{1}\mathrm{M}+\ldots+x_{p}\mathrm{M}))=i(\mathrm{M})\) (using the expression of \(i(\mathrm{M})\) given in a)).
\end{enumerate}

\begin{enumerate}
\def\labelenumi{\arabic{enumi})}
\setcounter{enumi}{8}
\tightlist
\item
  Let A be a Noetherian local ring.
\end{enumerate}

\begin{enumerate}
\def\labelenumi{\alph{enumi})}
\item
  Let M be a finitely generated Macaulay A-module, N a submodule of M containing \(\mathfrak{m}_{A}M\). Prove that N is a Buchsbaum module, and that one has \(i(N) = (\dim(M) - 1)[M/N : \kappa_{A}]\) .
\item
  If A is a Macaulay ring, the ideal \(\mathfrak{m}_{A}\) is a Buchsbaum module, which is not Macaulay if \(\dim(A) \geqslant 2\) .
\item
  Let \(k\) be a field; we take for A the subring of \(k[[X,Y]]\) generated by \(X^4, X^3Y, XY^3, Y^4\) . One has \(\dim(A) = 2\) , \(\text{prof}(A) = 1\) , and A is a Buchsbaum ring, with \(i(A) = 1\) (applying \(a\)) by taking for M the integral closure of A).
\end{enumerate}

¶ 10) a) Let \(\rho: A \to B\) be an injective homomorphism of Noetherian rings, such that the sub-A-module \(\mathrm{A}1_{B}\) of B is a direct factor. Prove that if B is a Macaulay ring, then so is A.

(Reduce to the case where A is local and, by reasoning by induction on \(\mathrm{prof}_{\mathrm{A}}(\mathrm{B})\) , where \(\mathrm{Ass}(\mathrm{B})\) contains a maximal ideal. Show, by considering a primary decomposition of 0 in B, that B is then isomorphic to a product, and conclude by reasoning by induction on the number of maximal ideals of B.

\begin{enumerate}
\def\labelenumi{\alph{enumi})}
\setcounter{enumi}{1}
\tightlist
\item
  Let B be a Macaulay ring, G a finite group of automorphisms of B whose order is invertible in B. Prove that the ring of invariants of G is a Macaulay ring.
\end{enumerate}

¶ 11) Let B be a Noetherian integrally closed ring, G a finite group of automorphisms of B, and A the ring of invariants of G.

\begin{enumerate}
\def\labelenumi{\alph{enumi})}
\tightlist
\item
  We assume that for every prime ideal p of \(\mathrm{A}\) of height 1 and every prime ideal q of B lying over p, the valuation \(v_{\mathfrak{q}}\) is unramified with respect to \(v_{\mathfrak{p}}\) (VI, § 8, no. 1). Let \(i: C(A) \to C(B)\) be the canonical homomorphism on divisor class groups (VII, § 1, no. 10); prove that the kernel of \(i\) is isomorphic to the cohomology group \(H^{1}(G,B^{*})\) (A, X, p. 111; let L be the fraction field of B; consider the exact sequence of Z \(^{(G)}\) -modules
\end{enumerate}

\[
0 \rightarrow \mathrm{B} ^ {*} \rightarrow \mathrm{L} ^ {*} \rightarrow \mathrm{D} (\mathrm{B}) \rightarrow \mathrm{C} (\mathrm{B}) \rightarrow 0,
\]

by observing that \(\mathrm{H}^{0}(\mathrm{G},\mathrm{D}(\mathrm{B}))\) is identified with \(\mathrm{D}(\mathrm{A})\) and that \(\mathrm{H}^{1}(\mathrm{G},\mathrm{L}^{*})\) is zero).

\begin{enumerate}
\def\labelenumi{\alph{enumi})}
\setcounter{enumi}{1}
\item
  Let k be a field of characteristic 2; take B to be the ring \(k[X_{1},\ldots,X_{4}]\) , and G to be the group of automorphisms of the k-algebra B generated by the automorphism \(\sigma\) such that \(\sigma(X_{i})=X_{i+1}\) for \(1\leqslant i\leqslant3\) and \(\sigma(X_{4})=X_{1}\) . Deduce from a) that the ring A of invariants of G is factorial.
\item
  Let m be the trace on \(\mathrm{A}\) of the ideal \((\mathrm{X}_1, \ldots, \mathrm{X}_4)\) ; set \(s = \mathrm{X}_1 + \ldots + \mathrm{X}_4\) , \(t = (\mathrm{X}_1 + \mathrm{X}_3)(\mathrm{X}_2 + \mathrm{X}_4)\) , \(u = (\mathrm{X}_1 + \mathrm{X}_2)(\mathrm{X}_3 + \mathrm{X}_4)(\mathrm{X}_1 + \mathrm{X}_4)(\mathrm{X}_2 + \mathrm{X}_3)\) , \(v = \mathrm{X}_1\mathrm{X}_2\mathrm{X}_3\mathrm{X}_4\) . Prove that the sequence \((s, t, u, v)\) is secant but not completely secant for \(\mathrm{A}_{\mathfrak{m}}\) . (Let \(x = \mathrm{X}_1^2(\mathrm{X}_2 + \mathrm{X}_3) + \mathrm{X}_2^2(\mathrm{X}_3 + \mathrm{X}_4) + \mathrm{X}_3^2(\mathrm{X}_4 + \mathrm{X}_1) + \mathrm{X}_4^2(\mathrm{X}_1 + \mathrm{X}_2)\) , \(y = \mathrm{X}_1\mathrm{X}_3 + \mathrm{X}_2\mathrm{X}_4\) , \(w = \mathrm{X}_1\mathrm{X}_2\mathrm{X}_3 + \mathrm{X}_2\mathrm{X}_3\mathrm{X}_4 + \mathrm{X}_3\mathrm{X}_4\mathrm{X}_1 + \mathrm{X}_4\mathrm{X}_1\mathrm{X}_2\) ; prove that we have \(sw + y^2 = u\) and that \(xy + tw\) is divisible by \(s\) , so that \(ux\) belongs to the ideal \((s, t)\) .)
\item
  Show that A is not a Macaulay ring, and that the sub-A-module \(A.1_{B}\) is not a direct summand of B.
\end{enumerate}

\exercisesection{Depth and homological dimension}

\begin{enumerate}
\def\labelenumi{\arabic{enumi})}
\tightlist
\item
  Let \(\rho: A \to B\) be a local homomorphism of Noetherian local rings, \(N\) a finitely generated B-module. Suppose that the A-modules \(B\) and \(N\) are flat.
\end{enumerate}

\begin{enumerate}
\def\labelenumi{\alph{enumi})}
\item
  Prove the equality \(\mathrm{dp}_{\mathrm{B}}(\mathrm{N}) = \mathrm{dp}_{\kappa_{\mathrm{A}} \otimes_{\mathrm{A}} \mathrm{B}}(\kappa_{\mathrm{A}} \otimes_{\mathrm{A}} \mathrm{N})\) (consider a resolution of N by finitely generated free B-modules).
\item
  Let M be a finitely generated A-module. Prove the equality
\end{enumerate}

\[
\mathrm{dp} _ {\mathrm{B}} (\mathrm{M} \otimes_ {\mathrm{A}} \mathrm{N}) = \mathrm{dp} _ {\mathrm{A}} (\mathrm{M}) + \mathrm{dp} _ {\mathrm{B}} (\mathrm{N}).
\]

(Let K (resp. L) be a resolution of the A-module M (resp. of the B-module N) by finitely generated free modules; prove that \(K \otimes_{A} L\) is a resolution of \(M \otimes_{A} N\) by free B-modules. If dp \(_{A}\) (M) = m and dp \(_{B}\) (N) = n, compute \(\mathrm{H}_{m+n}(\kappa_{\mathrm{B}} \otimes_{\mathrm{B}} (\mathrm{K} \otimes_{\mathrm{A}} \mathrm{L}))\) in terms of \(\mathrm{H}_{m}(\kappa_{\mathrm{A}} \otimes_{\mathrm{A}} \mathrm{K})\) and \(\mathrm{H}_{n}(\kappa_{\mathrm{B}} \otimes_{\mathrm{B}} \mathrm{L})\) .

\begin{enumerate}
\def\labelenumi{\arabic{enumi})}
\setcounter{enumi}{1}
\tightlist
\item
  Let A be a Noetherian ring, M an A-module admitting a finite free resolution (L, p) by finitely generated free modules.
\end{enumerate}

\begin{enumerate}
\def\labelenumi{\alph{enumi})}
\tightlist
\item
  Show that for every prime ideal p of A, one has
\end{enumerate}

\[
\sum_ {i \geqslant 0} (- 1) ^ {i} \operatorname{rg} _ {\mathrm{A}} (\mathrm{L} _ {i}) = \sum_ {i \geqslant 0} (- 1) ^ {i} \left[ \operatorname{Tor} _ {i} ^ {\mathrm{A}} (\mathrm{M}, \kappa (\mathfrak {p})): \kappa (\mathfrak {p}) \right] = \sum_ {i \geqslant 0} (- 1) ^ {i} \left[ \operatorname{Ext} _ {\mathrm{A}} ^ {i} (\mathrm{M}, \kappa (\mathfrak {p})): \kappa (\mathfrak {p}) \right];
\]

denote \(\chi(\mathrm{M})\) this integer.

\begin{enumerate}
\def\labelenumi{\alph{enumi})}
\setcounter{enumi}{1}
\item
  Prove that one has \(\chi(M) \geqslant 0\) (consider an ideal \(\mathfrak{p}\) associated to A).
\item
  Prove that the following conditions are equivalent:
\end{enumerate}

\begin{enumerate}
\def\labelenumi{(\roman{enumi})}
\item
  \(\chi (\mathbf{M}) = 0\)
\item
  the annihilator of M is not reduced to zero;
\item
  the annihilator of M contains a cancellable element.
\end{enumerate}

(If \(\chi(M) = 0\) , prove that \(M_{\mathfrak{p}}\) is zero for every \(\mathfrak{p} \in \operatorname{Ass}(A)\) ; if \(\chi(M) > 0\) , prove similarly that the annihilator of Ann(M) contains a cancellable element, which implies Ann(M) = 0.)

\begin{enumerate}
\def\labelenumi{\arabic{enumi})}
\setcounter{enumi}{2}
\tightlist
\item
  Let A be a Macaulay local ring. Recall (A, II, p.186, exerc. 16) that an ideal of A is said to be irreducible in A if it is distinct from A and is not the intersection of two ideals strictly containing it. Prove that the following conditions are equivalent:
\end{enumerate}

\begin{enumerate}
\def\labelenumi{(\roman{enumi})}
\item
  A is a Gorenstein ring;
\item
  every ideal of A generated by a maximal secant sequence is irreducible;
\item
  there exists a maximal secant sequence in A generating an irreducible ideal.
\end{enumerate}

(Reduce to the case where A has dimension 0, and consider the socle \(\mathrm{Hom}_{\mathrm{A}}(\kappa_{\mathrm{A}},\mathrm{A})\) of A.)

\begin{enumerate}
\def\labelenumi{\arabic{enumi})}
\setcounter{enumi}{3}
\tightlist
\item
  Let A be a Noetherian ring, M a nonzero finitely generated A-module. One has grade(M) ≤ dpA(M); M is said to be perfect if equality holds.
\end{enumerate}

\begin{enumerate}
\def\labelenumi{\alph{enumi})}
\tightlist
\item
  Let d be an integer; prove that the following conditions are equivalent:
\end{enumerate}

\begin{enumerate}
\def\labelenumi{(\roman{enumi})}
\item
  the A-module M is perfect of projective dimension d;
\item
  we have \(\operatorname{Ext}_{\mathrm{A}}^{p}(\mathrm{M},\mathrm{A}) = 0\) for \(p\neq d\)
\item
  for every prime (resp. maximal) ideal \(\mathfrak{p}\) in the support of M, the \(\mathrm{A}_{\mathfrak{p}}\) -module \(\mathrm{M}_{\mathfrak{p}}\) is perfect of projective dimension \(d\) .
\end{enumerate}

(First prove the equivalence of (i) and (ii).)

\begin{enumerate}
\def\labelenumi{\alph{enumi})}
\setcounter{enumi}{1}
\item
  Let J be an ideal of A generated by a completely secant sequence, and let m be an integer \(\geqslant 1\) . Prove that the A-module \(A/J^{m}\) is perfect.
\item
  If M is perfect, its associated prime ideals are the ideals p such that \(\text{prof}(A_{\mathfrak{p}}) = \text{dp}_{\mathrm{A}}(M)\) .
\item
  For every perfect module N of projective dimension d, set \(\mathrm{N}^{\prime} = \mathrm{Ext}_{\mathrm{A}}^{d}(\mathrm{N}, \mathrm{A})\) . Prove that \(N^{\prime}\) is perfect, of projective dimension d, and that \((\mathrm{N}^{\prime})^{\prime}\) is isomorphic to N (consider a projective resolution of N). We have \(\mathrm{Ass}(\mathrm{N}^{\prime}) = \mathrm{Ass}(\mathrm{N})\) .
\item
  Let \(x\) be a non-zero-divisor in the radical of A, such that the homothety \(x_{\mathrm{M}}\) is injective. For the A-module M to be perfect of projective dimension \(d\) , it is necessary and sufficient that the (A/xA)-module M/xM be so.
\item
  A Macaulay module of finite projective dimension is perfect (reduce to the local case, and argue by induction on dim(M)).
\item
  If A is a Macaulay ring, the perfect modules are the Macaulay modules of finite projective dimension.
\end{enumerate}

\begin{enumerate}
\def\labelenumi{\arabic{enumi})}
\setcounter{enumi}{4}
\tightlist
\item
  \begin{enumerate}
  \def\labelenumii{\alph{enumii})}
  \tightlist
  \item
    Let \(k\) be a field, and R a \(k\) -algebra of finite degree. For R to be a Gorenstein ring, it is necessary and sufficient that the R-module \(\mathrm{Hom}_{k}(\mathrm{R}, k)\) be free of rank 1 (reduce to the case where R is local, and observe that the R-module \(\mathrm{Hom}_{k}(\mathrm{R}, k)\) is injective and contains a submodule isomorphic to \(\kappa_{\mathrm{R}}\) ; for another proof, cf. § 8, no. 6, cor. of prop. 5).
  \end{enumerate}
\end{enumerate}

\begin{enumerate}
\def\labelenumi{\alph{enumi})}
\setcounter{enumi}{1}
\item
  Assume that the above conditions are satisfied, and moreover that the algebra R is graded of type N, with \(R_{0} = k\) . Let s be the largest integer such that \(R_{s} \neq 0\) . Prove that \(R_{s}\) is of dimension 1 over k, and equal to the socle of R; for every non-zero k-linear form \(\ell\) on \(R_{s}\) , the k-bilinear form \((x, y) \mapsto \ell(xy)\) on \(R_{i} \times R_{s- i}\) induces for every i an isomorphism of \(R_{i}\) on \(\operatorname{Hom}_{k}(R_{s-i}, k)\) .
\item
  Let A be a Noetherian ring, and B an A-algebra which is a flat module of finite type. Prove that the following conditions are equivalent:
\end{enumerate}

\begin{enumerate}
\def\labelenumi{(\roman{enumi})}
\item
  the B-module \(\operatorname{Hom}_{\mathrm{A}}(\mathbf{B},\mathbf{A})\) is projective of rank 1;
\item
  for every maximal (resp. prime) ideal \(\mathfrak{p}\) of A, the ring \(\kappa(\mathfrak{p})\otimes_{\mathrm{A}}\mathrm{B}\) is a Gorenstein ring.
\end{enumerate}

These conditions are satisfied in particular when B is a Gorenstein ring.

¶ 6) Let A be a Noetherian local ring, M a finitely generated A-module of finite projective dimension, and N an A-module such that \(\operatorname{Tor}_{i}^{\mathrm{A}}(\mathrm{M},\mathrm{N}) = 0\) for i \textgreater{} 0.

\begin{enumerate}
\def\labelenumi{\alph{enumi})}
\tightlist
\item
  Define for every integer n a canonical isomorphism
\end{enumerate}

\[
\operatorname{Ext} _ {\mathrm{A}} ^ {n} (\kappa_ {\mathrm{A}}, \mathrm{M} \otimes_ {\mathrm{A}} \mathrm{N}) \longrightarrow \bigoplus_ {p \geqslant n} \operatorname{Tor} _ {p - n} ^ {\mathrm{A}} (\kappa_ {\mathrm{A}}, \mathrm{M}) \otimes_ {k} \operatorname{Ext} _ {\mathrm{A}} ^ {p} (\kappa_ {\mathrm{A}}, \mathrm{N}),
\]

and prove that one has \(\mathrm{di}_{\mathrm{A}}(\mathrm{M}\otimes_{\mathrm{A}}\mathrm{N})\leqslant \mathrm{di}_{\mathrm{A}}(\mathrm{N})\).

(Let (P,p) be a finite free resolution of M, (I,e) an injective resolution of N, and C the complex \(P \otimes_A I\). Prove that the complex

\[
0 \to \mathrm{C} ^ {0} / \mathrm{B} ^ {0} (\mathrm{C}) \to \mathrm{C} ^ {1} \to \mathrm{C} ^ {2} \to \dots
\]

defines an injective resolution of \(M \otimes_A N\).)

\begin{enumerate}
\def\labelenumi{\alph{enumi})}
\setcounter{enumi}{1}
\item
  Prove the inequality \(\mathrm{prof}_{\mathrm{A}}(\mathrm{N}) = \mathrm{dp}_{\mathrm{A}}(\mathrm{M}) + \mathrm{prof}(\mathrm{M} \otimes_{\mathrm{A}} \mathrm{N})\) . Deduce another proof of the Auslander-Buchsbaum theorem.
\item
  Assume moreover that N is of finite projective dimension. Prove the equality \(\mathrm{dp}_{\mathrm{A}}(\mathrm{M}\otimes_{\mathrm{A}}\mathrm{N})=\mathrm{dp}_{\mathrm{A}}(\mathrm{M})+\mathrm{dp}_{\mathrm{A}}(\mathrm{N})\) .
\item
  Assume that A is a Gorenstein ring, and denote its dimension by d. Let n be an integer; show that the \(\kappa_{A}\) -vector spaces \(\operatorname{Ext}_{\mathrm{A}}^{n}(\mathbf{k}_{\mathrm{A}},\mathrm{M})\) , \(\operatorname{Tor}_{d-n}^{\mathrm{A}}(\mathbf{k}_{\mathrm{A}},\mathrm{M})\) and \(\operatorname{Ext}_{\mathrm{A}}^{d-n}(\mathrm{M},\mathbf{k}_{\mathrm{A}})\) have the same dimension.
\end{enumerate}

\begin{enumerate}
\def\labelenumi{\arabic{enumi})}
\setcounter{enumi}{6}
\tightlist
\item
  Let A be a Noetherian local ring, M a finitely generated A-module, and T a finitely generated A-module of finite injective dimension. Prove the equality
\end{enumerate}

\[
\operatorname{prof} (\mathrm{A}) = \operatorname{prof} _ {\mathrm{A}} (\mathrm{M}) + \sup \left\{i \mid \operatorname{Ext} _ {\mathrm{A}} ^ {i} (\mathrm{M}, \mathrm{T}) \neq 0 \right\}
\]

(reason by induction on \(\mathrm{prof}_{\mathrm{A}}(\mathrm{M})\) , using prop. 9).

\begin{enumerate}
\def\labelenumi{\arabic{enumi})}
\setcounter{enumi}{7}
\tightlist
\item
  Let A be a ring, and \(u: \mathrm{A}^p \to \mathrm{A}^q\) an A-linear map. We call the rank of \(u\) , and denote it by \(\text{rg}(u)\) , the largest integer \(r\) such that \(\wedge^r u \neq 0\) . One denotes by \(\mathfrak{d}(u)\) the ideal of A generated by the minors of order \(\text{rg}(u)\) of the matrix of \(u\) .
\end{enumerate}

\begin{enumerate}
\def\labelenumi{\alph{enumi})}
\item
  Show that for \(\mathfrak{d}(u)\) to be equal to A, it is necessary and sufficient that the cokernel of \(u\) be projective of rank \(q - \mathrm{rg}(u)\) (reduce to the case where \(\mathrm{A}\) is local and write the matrix of \(u\) in a suitable basis).
\item
  Let \(v: \mathbf{A}^q \to \mathrm{A}^s\) an A-linear map; we assume that we have \(\mathfrak{d}(u) = \mathfrak{d}(v) = \mathbf{A}\) For \(\operatorname{Ker} v = \operatorname{Im} u\) , it is necessary and sufficient that one has \(\operatorname{rg}(u) + \operatorname{rg}(v) = q\) (same method).
\item
  If \(\operatorname{rg}(u) = q\) , prove that \(\mathfrak{d}(u)\) annihilates the cokernel of \(u\) (by reducing to the case \(p = q\) ).
\end{enumerate}

Let A be a ring, and let (L, d) be a bounded complex of finitely generated free A-modules, zero in degrees \textless{} 0. We propose to prove that the following conditions are equivalent:

\begin{enumerate}
\def\labelenumi{(\roman{enumi})}
\item
  We have \(\mathrm{H}_i(\mathrm{L}) = 0\) for \(i > 0\) ;
\item
  We have \(\operatorname{prof}_{\mathrm{A}}(\mathfrak{d}(d_i); \mathrm{A}) \geqslant i\) and \(\operatorname{rg}(d_i) + \operatorname{rg}(d_{i+1}) = \operatorname{rg}_{\mathrm{A}}(\mathrm{L}_i)\) for every \(i \geqslant 1\) .
\end{enumerate}

\begin{enumerate}
\def\labelenumi{\alph{enumi})}
\item
  Show that it suffices to prove this statement when the ring A is Noetherian (consider the subring of A generated by the coefficients of the matrices \(d_{i}\) for \(i \geqslant 1\)). From now on, this condition is assumed to hold.
\item
  If A is a local ring of dimension 0, prove that (ii) implies (i) (use Exercise 8, b)).
\item
  Assume that the ring A is local, that condition (ii) is satisfied, and that \(\operatorname{Supp}(\mathrm{H}_i(\mathrm{L})) = \{\mathfrak{m}_{\mathrm{A}}\}\) for \(i \geqslant 1\) . Prove that \(\mathrm{H}_i(\mathrm{L}) = 0\) for \(i \geqslant 0\) (let \(p\) be the largest integer such that \(\mathfrak{d}(d_p) \neq \mathrm{A}\); deduce from exercise 8, b) that \(\mathrm{H}_i(\mathrm{L}) = 0\) for \(i > p\) , then from exerc. 8, a) and from cor. 2 of prop. 3 \(n^{\circ} 2\) ) that the complex \(0 \to L_p / B_p(L) \to L_{p-1} \to \ldots \to L_0\) is acyclic in degrees \(>0\) ).
\item
  Prove the implication (ii)⇒(i) (reduce to the case where A is local, and argue by induction on dim(A)).
\item
  We assume that the complex L is acyclic in degrees \textgreater{} 0. Let p be a prime ideal associated with A. Show that for every \(i \geqslant 1\) (Coker \(d_{i}\) )p is a free Ap-module, of rank \(\sum_{p \geqslant i-1} (-1)^{p} \operatorname{rg}_{\mathrm{A}}(\mathrm{L}_{p})\) ; moreover, we have \(\operatorname{rg}((d_{i})_{\mathfrak{p}}) = \sum_{p \geqslant i} (-1)^{p-i} \operatorname{rg}_{\mathrm{A}}(\mathrm{L}_{p})\) and \(\mathfrak{d}((d_{i})_{\mathfrak{p}}) = A_{\mathfrak{p}}\) (exerc. 8).
\end{enumerate}

\(f)\) Hence we have \(\operatorname{rg}(d_i) = \sum_{p \geqslant i} (-1)^{p - i} \operatorname{rg}_A(L_p)\) and consequently \(\operatorname{rg}(d_i) + \operatorname{rg}(d_{i+1}) = \operatorname{rg}_\mathrm{A}(L_i)\) for every \(i \geqslant 1\) .

\(g)\) Let \(i\) be an integer \(\geqslant 1\) , and \(\mathfrak{p}\) a prime ideal of \(\mathrm{A}\). Show that the condition \(\mathrm{dp}_{\mathrm{A}_{\mathfrak{p}}}(H_0(L)_{\mathfrak{p}}) \geqslant i\) is equivalent to \(\mathfrak{d}(d_i)A_{\mathfrak{p}} \neq A_{\mathfrak{p}}\) that is, \(\mathfrak{d}(d_i) \subset \mathfrak{p}\) (use exerc. 8, a)). Using prop. 8 of § 1, no. 5 and th. 1, deduce that one has \(\operatorname{prof}_{\mathrm{A}}(\mathfrak{d}(d_i); A) \geqslant i\) which proves the implication (i) \(\Rightarrow\) (ii).

¶ 10) We retain the hypotheses of Exerc. 9; we assume that the complex L is exact in degrees ≠ 0.

\begin{enumerate}
\def\labelenumi{\alph{enumi})}
\item
  Prove the inclusion \(V(\mathfrak{d}(d_{i+1})) \subset V(\mathfrak{d}(d_i))\) for every \(i \geqslant 1\) (observe that by Exercise 8, a prime ideal p does not contain \(\mathfrak{d}(d_i)\) if and only if the \(A_{\mathfrak{p}}\)-module \((\operatorname{Coker} d_i)_{\mathfrak{p}}\) is free).
\item
  We now assume that the annihilator of \(\mathrm{H}_{0}(\mathrm{L})\) is not reduced to zero. Prove that one has \(\mathrm{rg}(d_{1}) = \mathrm{rg}_{\mathrm{A}}(\mathrm{L}_{0})\) (if f is a cancellable element of \(\mathfrak{o}(d_{1})\) , the \(A_{f}\) -module \(\mathrm{H}_{0}(\mathrm{L})_{f}\) is projective of constant rank according to exerc. 8, hence necessarily zero); deduce from this that the support of \(\mathrm{H}_{0}(\mathrm{L})\) is \(\mathrm{V}(\mathfrak{o}(d_{1}))\) .
\item
  We set \(p = \text{prof}_A(\mathfrak{d}(d_1); A)\) . Prove that we have \(V(\mathfrak{d}(d_i)) = \text{Supp}(\mathbf{H}_0(L))\) for \(1 \leqslant i \leqslant p\) (reason as in the proof of \(a\) ), observing that if a prime ideal \(\mathfrak{p}\) of the support of \(\mathrm{H}_0(L)\) did not contain \(\mathfrak{d}(d_k)\) , one would have \(\text{dp}_{\mathrm{A}_{\mathfrak{p}}} \, H_0(L)_{\mathfrak{p}} < k \leqslant \text{grade}_{\mathrm{A}_{\mathfrak{p}}} \, H_0(L)_{\mathfrak{p}}\) .
\end{enumerate}

\begin{enumerate}
\def\labelenumi{\arabic{enumi})}
\setcounter{enumi}{10}
\tightlist
\item
  Let A be a local ring.
\end{enumerate}

\begin{enumerate}
\def\labelenumi{\alph{enumi})}
\tightlist
\item
  Let n be an integer \(\geqslant 1\) , \(u : A^{n-1} \to A^n\) an A-linear map, \(\mathfrak{d}(u)\) the ideal generated by the minors of order n-1 of u (exerc. 8). If \(\operatorname{prof}_{\mathrm{A}}(\mathfrak{d}(u); A) \geqslant 2\) , prove that \(\mathfrak{d}(u)\) admits the resolution
\end{enumerate}

\[
0 \rightarrow \mathrm{A} ^ {n - 1} \xrightarrow {u} \mathrm{A} ^ {n} \xrightarrow {v} \mathfrak {d} (u) \rightarrow 0,
\]

where \(v\) is induced by the homomorphism \(\wedge^{n-1}(^t u): A^n \to A\) (apply Exercise 9).

\begin{enumerate}
\def\labelenumi{\alph{enumi})}
\setcounter{enumi}{1}
\tightlist
\item
  Conversely, let J be an ideal of A of projective dimension 1. Prove that there exists a cancellable element a of A, an integer n, and an A-linear map \(u: A^{n-1} \to A^n\) such that \(\mathrm{J} = a\mathfrak{d}(u)\) ; one has \(\operatorname{prof}_{\mathrm{A}}(\mathrm{J}; \mathrm{A}) = 2\) ("Hilbert-Burch theorem": deduce from exerc. 9 that a minimal resolution of J is of the form \(0 \to A^{n-1} \xrightarrow{u} \mathrm{A}^n \to J\) , with \(\operatorname{prof}_{\mathrm{A}}(\mathfrak{d}(u); \mathrm{A}) \geqslant 2\) then apply a) and exerc. 1 of § 1, observing that \(\mathfrak{d}(u)\) contains a cancellable element of A).
\end{enumerate}

\begin{enumerate}
\def\labelenumi{\arabic{enumi})}
\setcounter{enumi}{11}
\tightlist
\item
  Let A be a Noetherian local ring, and let M and N be finitely generated non-zero \(\mathrm{A}\)-modules. Assume the projective dimension of M is finite. Let \(q\) be the largest integer such that \(\operatorname{Tor}_q^A(M, N) \neq 0\) .
\end{enumerate}

\begin{enumerate}
\def\labelenumi{\alph{enumi})}
\item
  Assume \(\mathrm{prof}_{\mathrm{A}}(\mathrm{Tor}_{q}^{\mathrm{A}}(\mathrm{M},\mathrm{N})) = 0\) . Prove the equality \(\mathrm{prof}_{\mathrm{A}}(\mathrm{N}) + q = \mathrm{dp}_{\mathrm{A}}(\mathrm{M})\) (one may reason by induction on \(\mathrm{prof}_{\mathrm{A}}(\mathrm{N})\) , considering an element \(x\) of \(\mathfrak{m}_{\mathrm{A}}\) a non-zero-divisor in N).
\item
  We further assume that the A-module \(M \otimes_{A} N\) is of finite length. Deduce from a) that \(\operatorname{Tor}_{i}^{A}(M, N)\) is zero for i \textgreater{} depth(A) - depth(M) - depth(N), and nonzero when equality holds. In particular, \(\operatorname{prof}(M) + \operatorname{prof}(N) \leqslant \operatorname{prof}(A)\) .
\end{enumerate}

¶ 13) Let A be a Noetherian local ring. Suppose the following condition is satisfied \(^{1}\) :

(GM) for every Noetherian local A-algebra B, there exists a B-module M such that \(\mathfrak{m}_{\mathrm{B}}\mathrm{M}\neq\mathrm{M}\) and \(\operatorname{prof}_{\mathrm{B}}(\mathrm{M})=\dim(\mathrm{B})\) .

\begin{enumerate}
\def\labelenumi{\alph{enumi})}
\item
  Let P be a complex of flat A-modules of finite length \(\ell\) such that \(\operatorname{Supp}(\mathrm{H}(\mathrm{P})) = \{\mathfrak{m}_{\mathrm{A}}\}\) . Prove that we have \(\ell \geqslant \dim(\mathrm{A})\) (use Exerc. 6 of § 1) \(^2\) .
\item
  Let B be a Noetherian ring, \(\rho: A \to B\) a homomorphism, \(^a\rho: \operatorname{Spec}(B) \to \operatorname{Spec}(A)\) the associated map, M a finitely generated A-module. Prove that the codimension of \(^a\rho^{-1}(\operatorname{Supp}(M))\) in \(\operatorname{Spec}(B)\) is \(\leqslant \mathrm{dp}_A(M)\) (applying \(a\) ) to the complex \(L \otimes_A B_q\) , where \(q\) is the minimal prime ideal of a component of \(^a\rho^{-1}(\operatorname{Supp}(M))\) and L a free resolution of M).
\item
  Let M, N be finitely generated A-modules, with \(M \neq 0\) . Prove the inequality
\end{enumerate}

\[
\dim (N) \leqslant \mathrm{dp} _ {A} (M) + \dim (M \otimes_ {A} N)
\]

(intersection theorem: reason by induction on the integer \(d = \dim(M \otimes_{\mathrm{A}} N)\) . If d = 0, apply b) to the ring B = A/Ann(M); in the general case consider an element of A that is regular for N and for M \(\otimes_{A} N\) ).

\begin{enumerate}
\def\labelenumi{\alph{enumi})}
\setcounter{enumi}{3}
\tightlist
\item
  For there to exist an A-module of finite length and finite projective dimension, it is necessary and sufficient that A be a Macaulay ring.
\end{enumerate}

\begin{enumerate}
\def\labelenumi{\arabic{enumi})}
\setcounter{enumi}{13}
\tightlist
\item
  Let \(\mathrm{A}\) be a Noetherian local ring satisfying condition (GM) (exerc. 13), M a finitely generated A-module of finite projective dimension.
\end{enumerate}

\begin{enumerate}
\def\labelenumi{\alph{enumi})}
\tightlist
\item
  Prove the inequalities
\end{enumerate}

\[
\dim (A) \leqslant \mathrm{dp} _ {A} (M) + \dim (M) \quad \dim (A) - \operatorname{prof} (A) \leqslant \dim (M) - \operatorname{prof} (M).
\]

If moreover the module M is perfect (exerc. 4), one has \(\dim(A) = \mathrm{dp}_{A}(M) + \dim(M)\) (use exerc. 15 of § 1).

\begin{enumerate}
\def\labelenumi{\alph{enumi})}
\setcounter{enumi}{1}
\item
  Let \(\mathfrak{p} \in \mathrm{Supp}(M)\) ; if the \(\mathrm{A}_{\mathfrak{p}}\) -module \(\mathrm{M}_{\mathfrak{p}}\) is Macaulay, the ring \(\mathrm{A}_{\mathfrak{p}}\) is a Macaulay ring. This is in particular the case if \(\mathfrak{p}\) is a minimal prime ideal of \(\mathrm{Supp}(M)\) .
\item
  Let \(\mathbf{x} = (x_{1},\ldots ,x_{r})\) be a regular sequence of elements of \(\mathfrak{m}_{\mathrm{A}}\) , which generates an ideal J of finite projective dimension. Prove that the sequence \(\mathbf{x}\) is completely regular for A (let \(\mathfrak{p}\in V(J)\) and let \(\mathfrak{q}\subset \mathfrak{p}\) be a minimal prime ideal of \(V(J)\) ; observe that one has \(\mathrm{prof}(A_{\mathfrak{p}})\geqslant \mathrm{dp}_{A_{\mathfrak{p}}}(A_{\mathfrak{p}} / J_{\mathfrak{p}})\geqslant \mathrm{dp}_{A_{\mathfrak{q}}}(A_{\mathfrak{q}} / J_{\mathfrak{q}})\) and \(\mathrm{dp}_{A_{\mathfrak{q}}}(A_{\mathfrak{q}} / J_{\mathfrak{q}})\geqslant r\) according to \(b\) ), whence finally \(\mathrm{prof}_{\mathrm{A}}(\mathrm{J};\mathrm{A})\geqslant r\) ; conclude that \(\mathbf{x}\) is completely regular with the aid of cor. 1 of th. 1 of § 1, no. 3).
\item
  Prove that every M-regular sequence is A-regular. (Prove that every element a of A whose annihilator is nonzero is not M-regular; for this, reduce by localization to the case where Supp(M) ∩ Supp(a) = \(\{\mathfrak{m}_{A}\}\) ; by applying exerc. 13 a) to the complex \(L \otimes_A A/p\), where L is a free resolution of M and p is an associated prime ideal of a, prove the inequalities \(\operatorname{prof}(A) \leqslant \dim(A/p) \leqslant \operatorname{dp}_{A}(M)\) and deduce from this that prof(M) = 0.)
\item
  For A to be an integral domain, it is necessary and sufficient that it contain a prime ideal of projective dimension \(< +\infty\) (if \(\mathrm{dp}_{\mathrm{A}}(\mathfrak{p}) < +\infty\) , show that A identifies with a subring of the regular ring \(\mathrm{A}_{\mathfrak{p}}\) ).
\end{enumerate}

\begin{enumerate}
\def\labelenumi{\arabic{enumi})}
\setcounter{enumi}{14}
\tightlist
\item
  Let A be a Noetherian local ring, M a finitely generated A-module of finite injective dimension.
\end{enumerate}

\begin{enumerate}
\def\labelenumi{\alph{enumi})}
\item
  Let \(\mathfrak{p} \in \operatorname{Supp}(M)\) . Prove the equality \(\operatorname{prof}(A) = \operatorname{prof}(A_{\mathfrak{p}}) + \dim(A/p)\) . (Let \(p = \operatorname{prof}(A_{\mathfrak{p}})\) and \(q = \dim(A/p)\) ; observe that one has \(p = \operatorname{di}_{A_{\mathfrak{p}}} (M_{\mathfrak{p}})\) , whence \(\operatorname{Ext}_{A_{\mathfrak{p}}}^{p}(\kappa(\mathfrak{p}), M_{\mathfrak{p}}) \neq 0\) and consequently \(\operatorname{Ext}_{A}^{p+q}(\kappa_A, M) \neq 0\) by lemma 3 of § 1, no. 7, which implies \(\operatorname{prof}(A) = \operatorname{di}_A(M) = p + q\) .) If \(\operatorname{Supp}(M) = \operatorname{Spec}(A)\), A is a Macaulay ring.
\item
  Prove the equality grade(M) + dim(M) = prof(A) (use a) and exercise 15 of § 1).
\end{enumerate}

¶ 16) Let A be a Macaulay ring and M an A-module of finite projective dimension. Prove that for M to be flat, it is necessary and sufficient that for every ideal J of A generated by a completely secant sequence, the canonical homomorphism \(J \otimes_{A} M \to JM\) is bijective (prove by induction on the length of the sequence that this condition implies \(\operatorname{Tor}_{i}^{A}(A/J, M) = 0\) for all i \textgreater{} 0, then, by descending induction on n, \(\operatorname{Tor}_{n}^{A}(N, M) = 0\) for every finitely generated A-module N and every integer n \textgreater{} 0). Example: a torsion-free module over a Dedekind ring is flat (cf. VII, § 2, exerc. 14).
